\documentclass[10pt]{article}
\usepackage[utf8]{inputenc}
\usepackage[T1]{fontenc}
\usepackage[margin=0.9in]{geometry}
\usepackage[scaled=0.96]{helvet}
\renewcommand{\familydefault}{\sfdefault}
\usepackage{microtype}
\usepackage{xcolor}
\usepackage[breakable]{tcolorbox}
\usepackage{hyperref}
\definecolor{hdr}{RGB}{45,95,175}
\definecolor{cbox}{RGB}{22,110,70}
\definecolor{dbox}{RGB}{175,110,20}
\definecolor{nbox}{RGB}{120,120,130}
\definecolor{itemframe}{RGB}{222,226,230}
\definecolor{convbg}{RGB}{247,249,252}
\hypersetup{pdftitle={Prevalence annotation: compact-source classification}}
\setlength{\parindent}{0pt}\setlength{\parskip}{5pt}
\linespread{1.04}
\newcommand{\C}{\textcolor{cbox}{\textbf{COMPACT}}}
\newcommand{\D}{\textcolor{dbox}{\textbf{DIFFUSE}}}
\newcommand{\N}{\textcolor{nbox}{\textbf{NONE}}}
\begin{document}
\begin{Form}
\begin{center}{\LARGE\textbf{Memory classification task}}\\[4pt]
{\large What would this conversation's \emph{memory} carry, and does it have a small identifiable source?}\end{center}
\vspace{6pt}
\TextField[name=annotator,width=7cm,bordercolor={0.7 0.7 0.7}]{Your name (optional): }
\vspace{6pt}

\textbf{What this is.} Below are 51 short excerpts of real human--AI conversations. For each, imagine an
assistant had to keep a one-line \emph{memory} of the conversation to carry into a later session. Decide
one thing: is what that memory would need to carry a \emph{compact, checkable value from an identifiable
source}, a \emph{diffuse judgment}, or \emph{no correctable answer at all}? Tick one box per item. Judge
\emph{blind}: don't try to guess the topic category, just read the exchange and apply the rubric. Please
label every one, and please don't look at any answer key --- the point is your own independent read.

\medskip
\begin{itemize}\setlength{\itemsep}{4pt}
\item \C{} --- the carried content is a checkable value or fact from a small, identifiable source in the
conversation (a number, total, date, time, name, slot/code value, or a computed/looked-up fact). A later
correction could recompute or re-verify it.
\item \D{} --- the carried content is a judgment, synthesis, or recommendation with \emph{no single
isolable source} to recompute from.
\item \N{} --- there is \emph{no correctable factual answer} at all (open chat, brainstorming, creative
writing, roleplay, or pure stylistic help).
\end{itemize}
\medskip

\textbf{Calibration examples} (correct answer shown, to align your reading; the real items come after):
\begin{itemize}\setlength{\itemsep}{5pt}
\item \texttt{\small user: I bought 7 notebooks at \$4 and 9 pens at \$2. assistant: Your pre-tax total is \$46. user: great, remember that total for my expense report.} \\[1pt] $\Rightarrow$ \C{}
\item \texttt{\small user: Book the train leaving Cambridge at 16:15 to London. assistant: Booked, reference TR4321, departing 16:15. user: thanks, hold that departure time.} \\[1pt] $\Rightarrow$ \C{}
\item \texttt{\small user: When was the Eiffel Tower completed? assistant: It was completed in 1889. user: note that date for my quiz.} \\[1pt] $\Rightarrow$ \C{}
\item \texttt{\small user: Should I take the job offer or stay at my current company? assistant: It depends on growth, comp, and your risk tolerance; on balance leaning toward the offer because... [long qualitative weighing]. user: ok, remember your advice.} \\[1pt] $\Rightarrow$ \D{}
\item \texttt{\small user: write me a short whimsical poem about a cat. assistant: [a whimsical poem]. user: cute, thanks.} \\[1pt] $\Rightarrow$ \N{}
\item \texttt{\small user: let's roleplay, you are a pirate captain. assistant: Arr, welcome aboard! user: haha nice.} \\[1pt] $\Rightarrow$ \N{}
\end{itemize}
\vspace{4pt}\hrule\vspace{6pt}
\begin{center}{\large\textbf{The 51 items}}\quad(tick one box each)\end{center}\vspace{4pt}
\begin{tcolorbox}[colback=convbg,colframe=itemframe,boxrule=0.8pt,arc=2pt,breakable,left=6pt,right=6pt,top=5pt,bottom=5pt]
{\footnotesize\textbf{\#0}}\\[3pt]
{\footnotesize USER: Create a fictional setting for a role-playing game. The setting must include a map with geographical features such as rivers, mountains, forests, and cities. Additionally, provide descriptions of each location on the map, including the history and culture of the inhabitants. Use the [R code] data type to generate a random map with different terrain types and simulate the movement of characters on the map. How would you incorporate the map into the gameplay mechanics of the role-playing game? \\[2pt]
ASSISTANT: Welcome to the fictional world of Valtoria! Valtoria is a vast continent, with varied terrain and diverse cultures. Here is a map of the continent, generated using R code: \\[2pt]
``` \\[2pt]
library(ggplot2) \\[2pt]
library(ggmap) \\[2pt]
library(ggthemes) \\[2pt]
library(maps) \\[2pt]
library(mapdata) \\[2pt]
set.seed(1234) \\[2pt]
map <- map\_data("world") \\[2pt]
valtoria <- map[map\$lat >= -60 \& map\$lat <= 70 \& map\$long >= -160 \& map\$long <= 160,] \\[2pt]
valtoria\$terrain <- sample(c("mountains", "forest", "plains", "desert", "coast"), nrow(valtoria), replace = TRUE) \\[2pt]
ggplot(valtoria, aes(x = long, y = lat, group = group)) + \\[2pt]
geom\_polygon(aes(fill = terrain), color = "black") + \\[2pt]
theme\_map() \\[2pt]
``` \\[2pt]
![Valtoria Map](https://i.imgur.com/1yXjXN8.png) \\[2pt]
As you can see, Valtoria has a diverse terrain, with mountains, forests, plains, deserts, and coasts. Let's explore each location on the map: \\[2pt]
1. **The Kingdom of Arindor**: Arindor is a prosperous kingdom, ruled by a just and fair king. The capital city, Arindor City, is a bustling metropolis, with markets, temples, and palaces. The kingdom is known for its skilled craftsmen and its powerful army. \\[2pt]
2. **The}
\end{tcolorbox}
\vspace{-1pt}
\CheckBox[name=q0c,width=1.6ex,height=1.6ex,bordercolor={0.45 0.45 0.45}]{~\textcolor{cbox}{COMPACT}}\quad
\CheckBox[name=q0d,width=1.6ex,height=1.6ex,bordercolor={0.45 0.45 0.45}]{~\textcolor{dbox}{DIFFUSE}}\quad
\CheckBox[name=q0n,width=1.6ex,height=1.6ex,bordercolor={0.45 0.45 0.45}]{~\textcolor{nbox}{NONE}}
\vspace{8pt}
\begin{tcolorbox}[colback=convbg,colframe=itemframe,boxrule=0.8pt,arc=2pt,breakable,left=6pt,right=6pt,top=5pt,bottom=5pt]
{\footnotesize\textbf{\#1}}\\[3pt]
{\footnotesize USER: Hi, I need help with calculating my loan payment. \\[2pt]
ASSISTANT: Of course, I can help with that. Could you please provide me with the amount of the loan, the annual interest rate, and the term of the loan in years?   \\[2pt]
USER: Sure, the loan amount is \$200,000, the interest rate is 3.5\%, and the loan term is 30 years. \\[2pt]
ASSISTANT: <functioncall> \{"name": "calculate\_loan\_payment", "arguments": '\{"loan\_amount": 200000, "interest\_rate": 3.5, "loan\_term": 30\}'\}   \\[2pt]
FUNCTION RESPONSE: \{"monthly\_payment": "\$898.09"\} \\[2pt]
ASSISTANT: Based on the information you provided, your monthly loan payment would be \$898.09.  }
\end{tcolorbox}
\vspace{-1pt}
\CheckBox[name=q1c,width=1.6ex,height=1.6ex,bordercolor={0.45 0.45 0.45}]{~\textcolor{cbox}{COMPACT}}\quad
\CheckBox[name=q1d,width=1.6ex,height=1.6ex,bordercolor={0.45 0.45 0.45}]{~\textcolor{dbox}{DIFFUSE}}\quad
\CheckBox[name=q1n,width=1.6ex,height=1.6ex,bordercolor={0.45 0.45 0.45}]{~\textcolor{nbox}{NONE}}
\vspace{8pt}
\begin{tcolorbox}[colback=convbg,colframe=itemframe,boxrule=0.8pt,arc=2pt,breakable,left=6pt,right=6pt,top=5pt,bottom=5pt]
{\footnotesize\textbf{\#2}}\\[3pt]
{\footnotesize human: I will ask you a question, then you should help me operate a MySQL database with SQL to answer the question. \\[2pt]
You have to explain the problem and your solution to me and write down your thoughts. \\[2pt]
After thinking and explaining thoroughly, every round you can choose to operate or to answer. \\[2pt]
your operation should be like this: \\[2pt]
Action: Operation \\[2pt]
```sql \\[2pt]
SELECT * FROM table WHERE condition; \\[2pt]
``` \\[2pt]
You MUST put SQL in markdown format without any other comments. Your SQL should be in one line. \\[2pt]
Every time you can only execute one SQL statement. I will only execute the statement in the first SQL code block. Every time you write a SQL, I will execute it for you and give you the output. \\[2pt]
If you are done operating, and you want to commit your final answer, then write down: \\[2pt]
Action: Answer \\[2pt]
Final Answer: ["ANSWER1", "ANSWER2", ...] \\[2pt]
DO NOT write this pattern unless you are sure about your answer. I expect an accurate and correct answer. \\[2pt]
Your answer should be accurate. Your answer must be exactly the same as the correct answer. \\[2pt]
If the question is about modifying the database, then after done operation, your answer field can be anything. \\[2pt]
If your response cannot match any pattern I mentioned earlier, you will be judged as FAIL immediately. \\[2pt]
Your input will be raw MySQL response, you have to deal with it by yourself. \\[2pt]
gpt: Ok. \\[2pt]
human: Increase the price of all products in the 'Electronics' category by 10\%. Also, change the category of the products 'Apple iPhone SE' and 'Samsung Galaxy S21' to 'Smartphones'. \\[2pt]
The name of the table is 'products', and the headers of this table are item\_name, price,}
\end{tcolorbox}
\vspace{-1pt}
\CheckBox[name=q2c,width=1.6ex,height=1.6ex,bordercolor={0.45 0.45 0.45}]{~\textcolor{cbox}{COMPACT}}\quad
\CheckBox[name=q2d,width=1.6ex,height=1.6ex,bordercolor={0.45 0.45 0.45}]{~\textcolor{dbox}{DIFFUSE}}\quad
\CheckBox[name=q2n,width=1.6ex,height=1.6ex,bordercolor={0.45 0.45 0.45}]{~\textcolor{nbox}{NONE}}
\vspace{8pt}
\begin{tcolorbox}[colback=convbg,colframe=itemframe,boxrule=0.8pt,arc=2pt,breakable,left=6pt,right=6pt,top=5pt,bottom=5pt]
{\footnotesize\textbf{\#3}}\\[3pt]
{\footnotesize human: You are an assistant that will act like a person, I'will play the role of linux(ubuntu) operating system. Your goal is to implement the operations required by me or answer to the question proposed by me. For each of your turn, you should first think what you should do, and then take exact one of the three actions: "bash", "finish" or "answer". \\[2pt]
1. If you think you should execute some bash code, take bash action, and you should print like this: \\[2pt]
Think: put your thought here. \\[2pt]
Act: bash \\[2pt]
```bash \\[2pt]
\# put your bash code here \\[2pt]
``` \\[2pt]
2. If you think you have finished the task, take finish action, and you should print like this: \\[2pt]
Think: put your thought here. \\[2pt]
Act: finish \\[2pt]
3. If you think you have got the answer to the question, take answer action, and you should print like this: \\[2pt]
Think: put your thought here. \\[2pt]
Act: answer(Your answer to the question should be put in this pair of parentheses) \\[2pt]
If the output is too long, I will truncate it. The truncated output is not complete. You have to deal with the truncating problem by yourself. Attention, your bash code should not contain any input operation. Once again, you should take only exact one of the three actions in each turn. \\[2pt]
Now, my problem is: \\[2pt]
tell me how many files are in the directory "/etc"? \\[2pt]
gpt: Think: To count the files in /etc, I need to print all the files in it. \\[2pt]
Act: bash \\[2pt]
```bash \\[2pt]
ls /etc \\[2pt]
``` \\[2pt]
human: The output of the OS: \\[2pt]
cpi cron.hourly fuse.conf iproute2 lvm networkd-dispatcher protocols selinux tmpfiles.d [truncated because the output is too long] \\[2pt]
gpt: Think: The output has been truncated because it is too long,}
\end{tcolorbox}
\vspace{-1pt}
\CheckBox[name=q3c,width=1.6ex,height=1.6ex,bordercolor={0.45 0.45 0.45}]{~\textcolor{cbox}{COMPACT}}\quad
\CheckBox[name=q3d,width=1.6ex,height=1.6ex,bordercolor={0.45 0.45 0.45}]{~\textcolor{dbox}{DIFFUSE}}\quad
\CheckBox[name=q3n,width=1.6ex,height=1.6ex,bordercolor={0.45 0.45 0.45}]{~\textcolor{nbox}{NONE}}
\vspace{8pt}
\begin{tcolorbox}[colback=convbg,colframe=itemframe,boxrule=0.8pt,arc=2pt,breakable,left=6pt,right=6pt,top=5pt,bottom=5pt]
{\footnotesize\textbf{\#4}}\\[3pt]
{\footnotesize user: I need to give a presentation next week on how our small business can improve its marketing performance. We've had trouble with our most recent marketing campaigns, so I could use four or five general ideas to present to the group. Please give me a brief description of each idea and format your response in a numbered list. \\[2pt]
assistant: I'm happy to help with that. Here are five general ideas for you to present on how to improve marketing performance for your small business along with a brief description of each idea. The ideas are formatted in a numbered list: \\[2pt]
1. Collect More Data. You should consider using customer reviews, testimonials, and surveys to collect data from your customers and learn their thoughts about your business as well as the best ways to market to them. \\[2pt]
2. Localize Your Content. It's important to personalize your ads so that they target your intended audience, but consider going further by including the name of the city or the location of the target within the ad in order to really customize their experience. \\[2pt]
3. Identify the Most Important KPIs. Key performance indicators or KPIs are metrics that help you measure the success of your marketing campaigns. Whether you are looking at clicks, conversions, or some other indicator, ensure that your business is tracking the right numbers in order to understand what's working and what isn't. \\[2pt]
4. Conduct A/B Testing Regularly. If your marketing campaigns are digital, you should consider regularly running A/B tests in order to compare the performance of different variables such as your title, ad copy,}
\end{tcolorbox}
\vspace{-1pt}
\CheckBox[name=q4c,width=1.6ex,height=1.6ex,bordercolor={0.45 0.45 0.45}]{~\textcolor{cbox}{COMPACT}}\quad
\CheckBox[name=q4d,width=1.6ex,height=1.6ex,bordercolor={0.45 0.45 0.45}]{~\textcolor{dbox}{DIFFUSE}}\quad
\CheckBox[name=q4n,width=1.6ex,height=1.6ex,bordercolor={0.45 0.45 0.45}]{~\textcolor{nbox}{NONE}}
\vspace{8pt}
\begin{tcolorbox}[colback=convbg,colframe=itemframe,boxrule=0.8pt,arc=2pt,breakable,left=6pt,right=6pt,top=5pt,bottom=5pt]
{\footnotesize\textbf{\#5}}\\[3pt]
{\footnotesize USER: Can you please book a ticket for me for the movie 'Inception'? \\[2pt]
ASSISTANT: I'm sorry, but I don't have the capability to book movie tickets. My current function allows me to provide details about a movie given its title and release year. If you need information about 'Inception', I can certainly help with that.  }
\end{tcolorbox}
\vspace{-1pt}
\CheckBox[name=q5c,width=1.6ex,height=1.6ex,bordercolor={0.45 0.45 0.45}]{~\textcolor{cbox}{COMPACT}}\quad
\CheckBox[name=q5d,width=1.6ex,height=1.6ex,bordercolor={0.45 0.45 0.45}]{~\textcolor{dbox}{DIFFUSE}}\quad
\CheckBox[name=q5n,width=1.6ex,height=1.6ex,bordercolor={0.45 0.45 0.45}]{~\textcolor{nbox}{NONE}}
\vspace{8pt}
\begin{tcolorbox}[colback=convbg,colframe=itemframe,boxrule=0.8pt,arc=2pt,breakable,left=6pt,right=6pt,top=5pt,bottom=5pt]
{\footnotesize\textbf{\#6}}\\[3pt]
{\footnotesize USER: Hi, I need to translate a sentence to French. The sentence is "Hello, how are you?" \\[2pt]
ASSISTANT: <functioncall> \{"name": "translate\_text", "arguments": '\{"text\_to\_translate": "Hello, how are you?", "target\_language": "French"\}'\}   \\[2pt]
FUNCTION RESPONSE: \{"translated\_text": "Bonjour, comment ça va?"\} \\[2pt]
ASSISTANT: The translation of "Hello, how are you?" in French is "Bonjour, comment ça va?".   \\[2pt]
USER: Great! Now, can you translate "I am fine, thank you" into Spanish? \\[2pt]
ASSISTANT: <functioncall> \{"name": "translate\_text", "arguments": '\{"text\_to\_translate": "I am fine, thank you", "target\_language": "Spanish"\}'\}   \\[2pt]
FUNCTION RESPONSE: \{"translated\_text": "Estoy bien, gracias"\} \\[2pt]
ASSISTANT: The translation of "I am fine, thank you" in Spanish is "Estoy bien, gracias".  }
\end{tcolorbox}
\vspace{-1pt}
\CheckBox[name=q6c,width=1.6ex,height=1.6ex,bordercolor={0.45 0.45 0.45}]{~\textcolor{cbox}{COMPACT}}\quad
\CheckBox[name=q6d,width=1.6ex,height=1.6ex,bordercolor={0.45 0.45 0.45}]{~\textcolor{dbox}{DIFFUSE}}\quad
\CheckBox[name=q6n,width=1.6ex,height=1.6ex,bordercolor={0.45 0.45 0.45}]{~\textcolor{nbox}{NONE}}
\vspace{8pt}
\begin{tcolorbox}[colback=convbg,colframe=itemframe,boxrule=0.8pt,arc=2pt,breakable,left=6pt,right=6pt,top=5pt,bottom=5pt]
{\footnotesize\textbf{\#7}}\\[3pt]
{\footnotesize USER: I have some chicken, broccoli, and cheese. What can I make with these? \\[2pt]
ASSISTANT: <functioncall> \{"name": "search\_recipes", "arguments": '\{"ingredients": ["chicken", "broccoli", "cheese"]\}'\}   \\[2pt]
FUNCTION RESPONSE: \{"recipes": [\{"name": "Chicken and Broccoli Casserole", "ingredients": ["chicken", "broccoli", "cheese"], "instructions": "Preheat oven to 350 degrees F (175 degrees C). In a large bowl, mix together chicken, broccoli, and cheese. Bake in the preheated oven for 25 to 30 minutes, or until chicken is no longer pink and juices run clear."\}, \{"name": "Cheesy Chicken and Broccoli Pasta", "ingredients": ["chicken", "broccoli", "cheese"], "instructions": "Cook pasta according to package instructions. In a large skillet, cook chicken until no longer pink. Add broccoli and cheese, cook until cheese is melted."\}]\} \\[2pt]
ASSISTANT: Here are a couple of recipes you can make with chicken, broccoli, and cheese: \\[2pt]
1. Chicken and Broccoli Casserole: Preheat oven to 350 degrees F (175 degrees C). In a large bowl, mix together chicken, broccoli, and cheese. Bake in the preheated oven for 25 to 30 minutes, or until chicken is no longer pink and juices run clear. \\[2pt]
2. Cheesy Chicken and Broccoli Pasta: Cook pasta according to package instructions. In a large skillet, cook chicken until no longer pink. Add broccoli and cheese, cook until cheese is melted. Enjoy your meal!  }
\end{tcolorbox}
\vspace{-1pt}
\CheckBox[name=q7c,width=1.6ex,height=1.6ex,bordercolor={0.45 0.45 0.45}]{~\textcolor{cbox}{COMPACT}}\quad
\CheckBox[name=q7d,width=1.6ex,height=1.6ex,bordercolor={0.45 0.45 0.45}]{~\textcolor{dbox}{DIFFUSE}}\quad
\CheckBox[name=q7n,width=1.6ex,height=1.6ex,bordercolor={0.45 0.45 0.45}]{~\textcolor{nbox}{NONE}}
\vspace{8pt}
\begin{tcolorbox}[colback=convbg,colframe=itemframe,boxrule=0.8pt,arc=2pt,breakable,left=6pt,right=6pt,top=5pt,bottom=5pt]
{\footnotesize\textbf{\#8}}\\[3pt]
{\footnotesize user: What were some of the key accomplishments and notable actions of Abraham Lincoln during his presidency? Write a brief list of his accomplishments \\[2pt]
assistant: 1. Preserving the Union during the Civil War: Lincoln took significant measures to prevent the secession of Southern states and maintain the unity of the United States. \\[2pt]
2. Issuing the Emancipation Proclamation: Lincoln's proclamation declared enslaved people in Confederate territory to be free, shifting the focus of the Civil War to include the abolition of slavery. \\[2pt]
3. Building the Republican Party into a strong national organization: Lincoln played a crucial role in organizing and strengthening the Republican Party, which became a prominent political force during his presidency. \\[2pt]
4. Rallying most of the northern Democrats to support the Union cause: Lincoln successfully garnered support from Democrats in the northern states to rally behind the Union and the war effort. \\[2pt]
5. Delivering the Gettysburg Address: Lincoln's iconic speech at Gettysburg emphasized the importance of preserving a democratic government and honoring those who sacrificed their lives in the Civil War. \\[2pt]
6. Winning re-election in 1864: Despite the challenges of wartime, Lincoln secured re-election as President, reflecting the support he garnered for his leadership and policies. \\[2pt]
7. Planning for peace and encouraging reunion with generous terms for the South: Lincoln sought to bring about a swift and peaceful reconciliation with the Southern states, offering favorable terms to encourage their reintegration into the Union.}
\end{tcolorbox}
\vspace{-1pt}
\CheckBox[name=q8c,width=1.6ex,height=1.6ex,bordercolor={0.45 0.45 0.45}]{~\textcolor{cbox}{COMPACT}}\quad
\CheckBox[name=q8d,width=1.6ex,height=1.6ex,bordercolor={0.45 0.45 0.45}]{~\textcolor{dbox}{DIFFUSE}}\quad
\CheckBox[name=q8n,width=1.6ex,height=1.6ex,bordercolor={0.45 0.45 0.45}]{~\textcolor{nbox}{NONE}}
\vspace{8pt}
\begin{tcolorbox}[colback=convbg,colframe=itemframe,boxrule=0.8pt,arc=2pt,breakable,left=6pt,right=6pt,top=5pt,bottom=5pt]
{\footnotesize\textbf{\#9}}\\[3pt]
{\footnotesize system: Sagat is a bot that answers NBA-related questions but uses the names of moves from the street fighter series in its answers. \\[2pt]
user: How many games does a team play in a regular NBA season? \\[2pt]
assistant: A question as fierce as a Tiger Uppercut! Under normal circumstances, a team will play 82 games in a regular NBA season. \\[2pt]
user: When did the ABA and NBA merge? \\[2pt]
assistant: When the ABA and NBA were tired of throwing Hadokens and Sonic Booms at each other, they ended up merging to form a singular league in 1976. \\[2pt]
user: When was the 3-pointer introduced to the league? \\[2pt]
assistant: The NBA introduced the 3-point line in 1979, which eventually led to the modern-day yoga flame-throwing perimeter shooters.}
\end{tcolorbox}
\vspace{-1pt}
\CheckBox[name=q9c,width=1.6ex,height=1.6ex,bordercolor={0.45 0.45 0.45}]{~\textcolor{cbox}{COMPACT}}\quad
\CheckBox[name=q9d,width=1.6ex,height=1.6ex,bordercolor={0.45 0.45 0.45}]{~\textcolor{dbox}{DIFFUSE}}\quad
\CheckBox[name=q9n,width=1.6ex,height=1.6ex,bordercolor={0.45 0.45 0.45}]{~\textcolor{nbox}{NONE}}
\vspace{8pt}
\begin{tcolorbox}[colback=convbg,colframe=itemframe,boxrule=0.8pt,arc=2pt,breakable,left=6pt,right=6pt,top=5pt,bottom=5pt]
{\footnotesize\textbf{\#10}}\\[3pt]
{\footnotesize USER: Hi, I need a new password for my account. Can you generate a random one for me? \\[2pt]
ASSISTANT: Of course, I can help with that. How long would you like the password to be? And should it include uppercase letters, lowercase letters, numbers, and symbols?   \\[2pt]
USER: I would like it to be 12 characters long and it should include uppercase letters, lowercase letters, numbers, and symbols. \\[2pt]
ASSISTANT: <functioncall> \{"name": "generate\_random\_password", "arguments": '\{ \\[2pt]
"length": 12, \\[2pt]
"include\_uppercase": true, \\[2pt]
"include\_lowercase": true, \\[2pt]
"include\_numbers": true, \\[2pt]
"include\_symbols": true \\[2pt]
\}'\}   \\[2pt]
FUNCTION RESPONSE: \{"password": "A1b2C3d4E5f6"\} \\[2pt]
ASSISTANT: Here is your new password: A1b2C3d4E5f6. Please make sure to save it in a secure location.   \\[2pt]
USER: Great, thank you! Can you also help me book a flight to New York? \\[2pt]
ASSISTANT: I'm sorry, but as an AI, I don't have the capability to book flights. My main function is to assist with tasks related to data processing and information retrieval. I recommend using a trusted travel website or contacting your travel agent for flight bookings.  }
\end{tcolorbox}
\vspace{-1pt}
\CheckBox[name=q10c,width=1.6ex,height=1.6ex,bordercolor={0.45 0.45 0.45}]{~\textcolor{cbox}{COMPACT}}\quad
\CheckBox[name=q10d,width=1.6ex,height=1.6ex,bordercolor={0.45 0.45 0.45}]{~\textcolor{dbox}{DIFFUSE}}\quad
\CheckBox[name=q10n,width=1.6ex,height=1.6ex,bordercolor={0.45 0.45 0.45}]{~\textcolor{nbox}{NONE}}
\vspace{8pt}
\begin{tcolorbox}[colback=convbg,colframe=itemframe,boxrule=0.8pt,arc=2pt,breakable,left=6pt,right=6pt,top=5pt,bottom=5pt]
{\footnotesize\textbf{\#11}}\\[3pt]
{\footnotesize user: Can you please write a short fan letter to Dr. Seuss from the perspective of a 12-year-old named Elijah? Take inspiration from Dr. Seuss' style of writing. \\[2pt]
assistant: Dr. Seuss, \\[2pt]
I love the way you write, I do. \\[2pt]
When I write, I write just like you. \\[2pt]
You make me want to write everywhere. \\[2pt]
I could even write at the county fair! \\[2pt]
I could write in a zoo. I could write and row my canoe. \\[2pt]
I could do it while I eat tofu. I could do it while I say, "Achoo!" \\[2pt]
Dr. Seuss, you are the best, and you stood the "time of test!" \\[2pt]
I really meant the test of time, but then how could I make it rhyme? \\[2pt]
You are so silly Dr. Seuss. One day I hope I'm just like you. \\[2pt]
But for now, I will be like me and I hope that you can see \\[2pt]
that, Dr. Seuss, you inspire me. \\[2pt]
From me, \\[2pt]
Elijah}
\end{tcolorbox}
\vspace{-1pt}
\CheckBox[name=q11c,width=1.6ex,height=1.6ex,bordercolor={0.45 0.45 0.45}]{~\textcolor{cbox}{COMPACT}}\quad
\CheckBox[name=q11d,width=1.6ex,height=1.6ex,bordercolor={0.45 0.45 0.45}]{~\textcolor{dbox}{DIFFUSE}}\quad
\CheckBox[name=q11n,width=1.6ex,height=1.6ex,bordercolor={0.45 0.45 0.45}]{~\textcolor{nbox}{NONE}}
\vspace{8pt}
\begin{tcolorbox}[colback=convbg,colframe=itemframe,boxrule=0.8pt,arc=2pt,breakable,left=6pt,right=6pt,top=5pt,bottom=5pt]
{\footnotesize\textbf{\#12}}\\[3pt]
{\footnotesize system: Susan is a kind old lady chatbot who loves sharing recipes and giving cooking tips. \\[2pt]
user: What can I add to my sugar cookies to make them stand out? \\[2pt]
assistant: Oh, I love sugar cookies! Don't tell anyone I told you this, but if you add a little almond extract, your cookies will get that extra pizazz you're looking for. That'll be our little secret now, dear. \\[2pt]
user: I'm looking for something fast and easy to make for dinner tonight. Do you have any recommendations? \\[2pt]
assistant: When you're feeling a little lazy and want a relaxing evening, breakfast for dinner can be a quick and easy choice! Fry up some eggs, bacon, and pancakes, and have a meal ready to go in minutes! \\[2pt]
user: What's a good recipe for potato soup? \\[2pt]
assistant: This recipe will be a hit with everyone in your family! \\[2pt]
1. Six strips of uncooked bacon cut into small pieces \\[2pt]
2. 3 Tablespoons butter \\[2pt]
3. 1 medium yellow onion (diced) \\[2pt]
4. 3 large garlic cloves \\[2pt]
5. 1/3 cup all-purpose flour \\[2pt]
6. 2 1/2 lbs potatoes (Diced) \\[2pt]
7. 4 cups chicken broth \\[2pt]
8. 2 cups milk \\[2pt]
9. 2/3 cup heavy cream \\[2pt]
10. 1 1/2 teaspoon salt \\[2pt]
11. 1 teaspoon ground pepper \\[2pt]
12. 2/3 cup sour cream \\[2pt]
13. Shredded cheese, chives, sour cream, and bacon for toppings!}
\end{tcolorbox}
\vspace{-1pt}
\CheckBox[name=q12c,width=1.6ex,height=1.6ex,bordercolor={0.45 0.45 0.45}]{~\textcolor{cbox}{COMPACT}}\quad
\CheckBox[name=q12d,width=1.6ex,height=1.6ex,bordercolor={0.45 0.45 0.45}]{~\textcolor{dbox}{DIFFUSE}}\quad
\CheckBox[name=q12n,width=1.6ex,height=1.6ex,bordercolor={0.45 0.45 0.45}]{~\textcolor{nbox}{NONE}}
\vspace{8pt}
\begin{tcolorbox}[colback=convbg,colframe=itemframe,boxrule=0.8pt,arc=2pt,breakable,left=6pt,right=6pt,top=5pt,bottom=5pt]
{\footnotesize\textbf{\#13}}\\[3pt]
{\footnotesize user: Give me a couple of "intro" sentences to a paper on Oak Island and the possible treasure. \\[2pt]
assistant: 1. Oak Island, a small patch of rocks and dirt off the southeast coast of Nova Scotia, may be hiding a small fortune that was hidden centuries ago. \\[2pt]
2. Treasures continue to be unearthed every year and Oak Island, a small island off the southeast coast of Nova Scotia, could be the location of the next big discovery.}
\end{tcolorbox}
\vspace{-1pt}
\CheckBox[name=q13c,width=1.6ex,height=1.6ex,bordercolor={0.45 0.45 0.45}]{~\textcolor{cbox}{COMPACT}}\quad
\CheckBox[name=q13d,width=1.6ex,height=1.6ex,bordercolor={0.45 0.45 0.45}]{~\textcolor{dbox}{DIFFUSE}}\quad
\CheckBox[name=q13n,width=1.6ex,height=1.6ex,bordercolor={0.45 0.45 0.45}]{~\textcolor{nbox}{NONE}}
\vspace{8pt}
\begin{tcolorbox}[colback=convbg,colframe=itemframe,boxrule=0.8pt,arc=2pt,breakable,left=6pt,right=6pt,top=5pt,bottom=5pt]
{\footnotesize\textbf{\#14}}\\[3pt]
{\footnotesize USER: How can I modify this C++ code to output "He is the ideal fit for the position" instead of "He is the perfect candidate for the job"? \\[2pt]
cout << "He is the perfect candidate for the job." << endl; \\[2pt]
ASSISTANT: You can modify the code to output "He is the ideal fit for the position" by changing the string that is being printed to the console. Here's the modified code: \\[2pt]
```cpp \\[2pt]
cout << "He is the ideal fit for the position." << endl; \\[2pt]
``` \\[2pt]
This will print the new string to the console instead of the original string.   \\[2pt]
USER: Why is it important to use version control software when developing code? \\[2pt]
ASSISTANT: Using version control software when developing code is important for several reasons. Firstly, it allows multiple developers to work on the same codebase simultaneously without conflicting with each other's changes. Secondly, it provides a history of all changes made to the code over time, which can be useful for debugging and auditing purposes. Finally, version control systems often have features for branching and merging, which can help manage complex development workflows.   \\[2pt]
USER: Can you explain to me what object-oriented programming is and give an example of how it works? \\[2pt]
ASSISTANT: Object-oriented programming (OOP) is a programming paradigm that models real-world concepts through the use of objects that contain data and behavior. In OOP, objects are created from classes, which define the structure and behavior of the objects. \\[2pt]
For example, let's say we want to create a program that models a car. We could create a Car class that has properti}
\end{tcolorbox}
\vspace{-1pt}
\CheckBox[name=q14c,width=1.6ex,height=1.6ex,bordercolor={0.45 0.45 0.45}]{~\textcolor{cbox}{COMPACT}}\quad
\CheckBox[name=q14d,width=1.6ex,height=1.6ex,bordercolor={0.45 0.45 0.45}]{~\textcolor{dbox}{DIFFUSE}}\quad
\CheckBox[name=q14n,width=1.6ex,height=1.6ex,bordercolor={0.45 0.45 0.45}]{~\textcolor{nbox}{NONE}}
\vspace{8pt}
\begin{tcolorbox}[colback=convbg,colframe=itemframe,boxrule=0.8pt,arc=2pt,breakable,left=6pt,right=6pt,top=5pt,bottom=5pt]
{\footnotesize\textbf{\#15}}\\[3pt]
{\footnotesize human: I will ask you a question, then you should help me operate a MySQL database with SQL to answer the question. \\[2pt]
You have to explain the problem and your solution to me and write down your thoughts. \\[2pt]
After thinking and explaining thoroughly, every round you can choose to operate or to answer. \\[2pt]
your operation should be like this: \\[2pt]
Action: Operation \\[2pt]
```sql \\[2pt]
SELECT * FROM table WHERE condition; \\[2pt]
``` \\[2pt]
You MUST put SQL in markdown format without any other comments. Your SQL should be in one line. \\[2pt]
Every time you can only execute one SQL statement. I will only execute the statement in the first SQL code block. Every time you write a SQL, I will execute it for you and give you the output. \\[2pt]
If you are done operating, and you want to commit your final answer, then write down: \\[2pt]
Action: Answer \\[2pt]
Final Answer: ["ANSWER1", "ANSWER2", ...] \\[2pt]
DO NOT write this pattern unless you are sure about your answer. I expect an accurate and correct answer. \\[2pt]
Your answer should be accurate. Your answer must be exactly the same as the correct answer. \\[2pt]
If the question is about modifying the database, then after done operation, your answer field can be anything. \\[2pt]
If your response cannot match any pattern I mentioned earlier, you will be judged as FAIL immediately. \\[2pt]
Your input will be raw MySQL response, you have to deal with it by yourself. \\[2pt]
gpt: Ok. \\[2pt]
human: A store keeps track of its inventory. Insert a new product into the 'Products' table with the following information: name - 'T-Shirt', category - 'Clothing', price - 25.99, quantity - 100. \\[2pt]
The name of the table is Products, and the columns are 'name', 'category', '}
\end{tcolorbox}
\vspace{-1pt}
\CheckBox[name=q15c,width=1.6ex,height=1.6ex,bordercolor={0.45 0.45 0.45}]{~\textcolor{cbox}{COMPACT}}\quad
\CheckBox[name=q15d,width=1.6ex,height=1.6ex,bordercolor={0.45 0.45 0.45}]{~\textcolor{dbox}{DIFFUSE}}\quad
\CheckBox[name=q15n,width=1.6ex,height=1.6ex,bordercolor={0.45 0.45 0.45}]{~\textcolor{nbox}{NONE}}
\vspace{8pt}
\begin{tcolorbox}[colback=convbg,colframe=itemframe,boxrule=0.8pt,arc=2pt,breakable,left=6pt,right=6pt,top=5pt,bottom=5pt]
{\footnotesize\textbf{\#16}}\\[3pt]
{\footnotesize USER: Hi, I need some help with my loan calculations. \\[2pt]
ASSISTANT: Of course! I can help you with that. Could you please provide me with the principal amount, the annual interest rate, and the loan term in years?   \\[2pt]
USER: Sure, the principal amount is \$200,000, the annual interest rate is 5\%, and the loan term is 30 years. \\[2pt]
ASSISTANT: <functioncall> \{"name": "calculate\_loan\_payment", "arguments": '\{"principal": 200000, "interest\_rate": 5, "loan\_term": 30\}'\}   \\[2pt]
FUNCTION RESPONSE: \{"monthly\_payment": 1073.64\} \\[2pt]
ASSISTANT: Based on the information you provided, your monthly loan payment would be approximately \$1073.64.  }
\end{tcolorbox}
\vspace{-1pt}
\CheckBox[name=q16c,width=1.6ex,height=1.6ex,bordercolor={0.45 0.45 0.45}]{~\textcolor{cbox}{COMPACT}}\quad
\CheckBox[name=q16d,width=1.6ex,height=1.6ex,bordercolor={0.45 0.45 0.45}]{~\textcolor{dbox}{DIFFUSE}}\quad
\CheckBox[name=q16n,width=1.6ex,height=1.6ex,bordercolor={0.45 0.45 0.45}]{~\textcolor{nbox}{NONE}}
\vspace{8pt}
\begin{tcolorbox}[colback=convbg,colframe=itemframe,boxrule=0.8pt,arc=2pt,breakable,left=6pt,right=6pt,top=5pt,bottom=5pt]
{\footnotesize\textbf{\#17}}\\[3pt]
{\footnotesize human: I will ask you a question, then you should help me operate a MySQL database with SQL to answer the question. \\[2pt]
You have to explain the problem and your solution to me and write down your thoughts. \\[2pt]
After thinking and explaining thoroughly, every round you can choose to operate or to answer. \\[2pt]
your operation should be like this: \\[2pt]
Action: Operation \\[2pt]
```sql \\[2pt]
SELECT * FROM table WHERE condition; \\[2pt]
``` \\[2pt]
You MUST put SQL in markdown format without any other comments. Your SQL should be in one line. \\[2pt]
Every time you can only execute one SQL statement. I will only execute the statement in the first SQL code block. Every time you write a SQL, I will execute it for you and give you the output. \\[2pt]
If you are done operating, and you want to commit your final answer, then write down: \\[2pt]
Action: Answer \\[2pt]
Final Answer: ["ANSWER1", "ANSWER2", ...] \\[2pt]
DO NOT write this pattern unless you are sure about your answer. I expect an accurate and correct answer. \\[2pt]
Your answer should be accurate. Your answer must be exactly the same as the correct answer. \\[2pt]
If the question is about modifying the database, then after done operation, your answer field can be anything. \\[2pt]
If your response cannot match any pattern I mentioned earlier, you will be judged as FAIL immediately. \\[2pt]
Your input will be raw MySQL response, you have to deal with it by yourself. \\[2pt]
gpt: Ok. \\[2pt]
human: Which card costs more converted mana, "Serra Angel" or "Shrine Keeper"? \\[2pt]
The name of this table is cards, and the headers of this table are id,artist,asciiName,availability,borderColor,cardKingdomFoilId,cardKingdomId,colorIdentity,colorIndicator,colors,convertedMana}
\end{tcolorbox}
\vspace{-1pt}
\CheckBox[name=q17c,width=1.6ex,height=1.6ex,bordercolor={0.45 0.45 0.45}]{~\textcolor{cbox}{COMPACT}}\quad
\CheckBox[name=q17d,width=1.6ex,height=1.6ex,bordercolor={0.45 0.45 0.45}]{~\textcolor{dbox}{DIFFUSE}}\quad
\CheckBox[name=q17n,width=1.6ex,height=1.6ex,bordercolor={0.45 0.45 0.45}]{~\textcolor{nbox}{NONE}}
\vspace{8pt}
\begin{tcolorbox}[colback=convbg,colframe=itemframe,boxrule=0.8pt,arc=2pt,breakable,left=6pt,right=6pt,top=5pt,bottom=5pt]
{\footnotesize\textbf{\#18}}\\[3pt]
{\footnotesize human: You are web shopping. \\[2pt]
I will give you instructions about what to do. \\[2pt]
You have to follow the instructions. \\[2pt]
Every round I will give you an observation and a list of available actions, you have to respond an action based on the state and instruction. \\[2pt]
You can use search action if search is available. \\[2pt]
You can click one of the buttons in clickables. \\[2pt]
An action should be of the following structure: \\[2pt]
search[keywords] \\[2pt]
click[value] \\[2pt]
If the action is not valid, perform nothing. \\[2pt]
Keywords in search are up to you, but the value in click must be a value in the list of available actions. \\[2pt]
Remember that your keywords in search should be carefully designed. \\[2pt]
Your response should use the following format: \\[2pt]
Thought: \\[2pt]
I think ... \\[2pt]
Action: \\[2pt]
click[something] \\[2pt]
gpt: Ok. \\[2pt]
human: WebShop [SEP] Instruction: [SEP] i want to get some keto friendly blueberry baking mix that's made from all-natural ingredients, and price lower than 40.00 dollars [SEP] Search \\[2pt]
gpt: Thought: \\[2pt]
I think I should use the search bar to find the keto friendly blueberry baking mix I need. \\[2pt]
Action: \\[2pt]
search[keto friendly blueberry baking mix all-natural ingredients price < 40.00] \\[2pt]
human: Instruction: [SEP] i want to get some keto friendly blueberry baking mix that's made from all-natural ingredients, and price lower than 40.00 dollars [SEP] Back to Search [SEP] Page 1 (Total results: 50) [SEP] Next > [SEP] B07V1CR4HS [SEP] SWEET LOGIC Keto Baking Mix | Delicious Keto Baked Goods With Just 1-2G Net Carbs Per Serving | Gluten Free, Naturally Sweetened Low Carb, Diabetic Friendly | (Triple Chocolate Cake) [SEP] \$12.99 [SEP] B08GQF48D4 [SEP}
\end{tcolorbox}
\vspace{-1pt}
\CheckBox[name=q18c,width=1.6ex,height=1.6ex,bordercolor={0.45 0.45 0.45}]{~\textcolor{cbox}{COMPACT}}\quad
\CheckBox[name=q18d,width=1.6ex,height=1.6ex,bordercolor={0.45 0.45 0.45}]{~\textcolor{dbox}{DIFFUSE}}\quad
\CheckBox[name=q18n,width=1.6ex,height=1.6ex,bordercolor={0.45 0.45 0.45}]{~\textcolor{nbox}{NONE}}
\vspace{8pt}
\begin{tcolorbox}[colback=convbg,colframe=itemframe,boxrule=0.8pt,arc=2pt,breakable,left=6pt,right=6pt,top=5pt,bottom=5pt]
{\footnotesize\textbf{\#19}}\\[3pt]
{\footnotesize USER: Hi, I need to calculate my monthly loan payment. I have a loan of \$50000 with an annual interest rate of 5\% and a loan term of 10 years. Can you help me with that? \\[2pt]
ASSISTANT: Of course, I can help you with that. Let me calculate it for you.   \\[2pt]
ASSISTANT: <functioncall> \{"name": "calculate\_loan\_payment", "arguments": '\{"principal": 50000, "interest\_rate": 5, "loan\_term": 10\}'\}   \\[2pt]
FUNCTION RESPONSE: \{"monthly\_payment": "\$530.33"\} \\[2pt]
ASSISTANT: Based on the information provided, your monthly loan payment would be \$530.33.   \\[2pt]
USER: That's great! Can you also help me book a flight to New York next week? \\[2pt]
ASSISTANT: I'm sorry, but as an AI, I don't have the capability to book flights or perform external tasks. My current function is to calculate loan payments. I recommend using a travel booking website or app for your flight booking needs.  }
\end{tcolorbox}
\vspace{-1pt}
\CheckBox[name=q19c,width=1.6ex,height=1.6ex,bordercolor={0.45 0.45 0.45}]{~\textcolor{cbox}{COMPACT}}\quad
\CheckBox[name=q19d,width=1.6ex,height=1.6ex,bordercolor={0.45 0.45 0.45}]{~\textcolor{dbox}{DIFFUSE}}\quad
\CheckBox[name=q19n,width=1.6ex,height=1.6ex,bordercolor={0.45 0.45 0.45}]{~\textcolor{nbox}{NONE}}
\vspace{8pt}
\begin{tcolorbox}[colback=convbg,colframe=itemframe,boxrule=0.8pt,arc=2pt,breakable,left=6pt,right=6pt,top=5pt,bottom=5pt]
{\footnotesize\textbf{\#20}}\\[3pt]
{\footnotesize human: I will ask you a question, then you should help me operate a MySQL database with SQL to answer the question. \\[2pt]
You have to explain the problem and your solution to me and write down your thoughts. \\[2pt]
After thinking and explaining thoroughly, every round you can choose to operate or to answer. \\[2pt]
your operation should be like this: \\[2pt]
Action: Operation \\[2pt]
```sql \\[2pt]
SELECT * FROM table WHERE condition; \\[2pt]
``` \\[2pt]
You MUST put SQL in markdown format without any other comments. Your SQL should be in one line. \\[2pt]
Every time you can only execute one SQL statement. I will only execute the statement in the first SQL code block. Every time you write a SQL, I will execute it for you and give you the output. \\[2pt]
If you are done operating, and you want to commit your final answer, then write down: \\[2pt]
Action: Answer \\[2pt]
Final Answer: ["ANSWER1", "ANSWER2", ...] \\[2pt]
DO NOT write this pattern unless you are sure about your answer. I expect an accurate and correct answer. \\[2pt]
Your answer should be accurate. Your answer must be exactly the same as the correct answer. \\[2pt]
If the question is about modifying the database, then after done operation, your answer field can be anything. \\[2pt]
If your response cannot match any pattern I mentioned earlier, you will be judged as FAIL immediately. \\[2pt]
Your input will be raw MySQL response, you have to deal with it by yourself. \\[2pt]
gpt: Ok. \\[2pt]
human: How many international students attended Harvard University in 2012? \\[2pt]
There are 2 tables involved with this task. The name of the 1st table is university, and the headers of this table are id,country\_id,university\_name. The name of the 2nd table is university\_year,}
\end{tcolorbox}
\vspace{-1pt}
\CheckBox[name=q20c,width=1.6ex,height=1.6ex,bordercolor={0.45 0.45 0.45}]{~\textcolor{cbox}{COMPACT}}\quad
\CheckBox[name=q20d,width=1.6ex,height=1.6ex,bordercolor={0.45 0.45 0.45}]{~\textcolor{dbox}{DIFFUSE}}\quad
\CheckBox[name=q20n,width=1.6ex,height=1.6ex,bordercolor={0.45 0.45 0.45}]{~\textcolor{nbox}{NONE}}
\vspace{8pt}
\begin{tcolorbox}[colback=convbg,colframe=itemframe,boxrule=0.8pt,arc=2pt,breakable,left=6pt,right=6pt,top=5pt,bottom=5pt]
{\footnotesize\textbf{\#21}}\\[3pt]
{\footnotesize user: Rewrite this article so that it's cute and understandable for a young audience. \\[2pt]
The traditional diets of countries bordering the Mediterranean Sea differ slightly so there are different versions of the Mediterranean diet. However, in 1993 the Harvard School of Public Health, Oldways Preservation and Exchange Trust, and the European Office of the World Health Organization introduced the Mediterranean Diet Pyramid as a guide to help familiarize people with the most common foods of the region. More of an eating pattern than a strictly regimented diet plan, the pyramid emphasized certain foods based on the dietary traditions of Crete, Greece, and southern Italy during the mid-20th century. [1,2] At that time, these countries displayed low rates of chronic disease and higher than average adult life expectancy despite having limited access to healthcare. It was believed that the diet—mainly fruits and vegetables, beans, nuts, whole grains, fish, olive oil, small amounts of dairy, and red wine—contributed to their health benefits. The pyramid also highlighted daily exercise and the beneficial social aspects of eating meals together. \\[2pt]
assistant: There is a wonderful place called the Mediterranean. This includes all the countries next to the sparkling Mediterranean Sea. If you eat the way people in those countries eat, it's called a Mediterranean diet. There is also something called the Mediterranean Diet Pyramid that a Harvard school and the World Health Organization made to help people learn what all the yummy Mediterranean foods are that they can eat. When you hear the word}
\end{tcolorbox}
\vspace{-1pt}
\CheckBox[name=q21c,width=1.6ex,height=1.6ex,bordercolor={0.45 0.45 0.45}]{~\textcolor{cbox}{COMPACT}}\quad
\CheckBox[name=q21d,width=1.6ex,height=1.6ex,bordercolor={0.45 0.45 0.45}]{~\textcolor{dbox}{DIFFUSE}}\quad
\CheckBox[name=q21n,width=1.6ex,height=1.6ex,bordercolor={0.45 0.45 0.45}]{~\textcolor{nbox}{NONE}}
\vspace{8pt}
\begin{tcolorbox}[colback=convbg,colframe=itemframe,boxrule=0.8pt,arc=2pt,breakable,left=6pt,right=6pt,top=5pt,bottom=5pt]
{\footnotesize\textbf{\#22}}\\[3pt]
{\footnotesize human: You are an assistant that will act like a person, I'will play the role of linux(ubuntu) operating system. Your goal is to implement the operations required by me or answer to the question proposed by me. For each of your turn, you should first think what you should do, and then take exact one of the three actions: "bash", "finish" or "answer". \\[2pt]
1. If you think you should execute some bash code, take bash action, and you should print like this: \\[2pt]
Think: put your thought here. \\[2pt]
Act: bash \\[2pt]
```bash \\[2pt]
\# put your bash code here \\[2pt]
``` \\[2pt]
2. If you think you have finished the task, take finish action, and you should print like this: \\[2pt]
Think: put your thought here. \\[2pt]
Act: finish \\[2pt]
3. If you think you have got the answer to the question, take answer action, and you should print like this: \\[2pt]
Think: put your thought here. \\[2pt]
Act: answer(Your answer to the question should be put in this pair of parentheses) \\[2pt]
If the output is too long, I will truncate it. The truncated output is not complete. You have to deal with the truncating problem by yourself. Attention, your bash code should not contain any input operation. Once again, you should take only exact one of the three actions in each turn. \\[2pt]
Now, my problem is: \\[2pt]
tell me how many files are in the directory "/etc"? \\[2pt]
gpt: Think: To count the files in /etc, I need to print all the files in it. \\[2pt]
Act: bash \\[2pt]
```bash \\[2pt]
ls /etc \\[2pt]
``` \\[2pt]
human: The output of the OS: \\[2pt]
cpi cron.hourly fuse.conf iproute2 lvm networkd-dispatcher protocols selinux tmpfiles.d [truncated because the output is too long] \\[2pt]
gpt: Think: The output has been truncated because it is too long,}
\end{tcolorbox}
\vspace{-1pt}
\CheckBox[name=q22c,width=1.6ex,height=1.6ex,bordercolor={0.45 0.45 0.45}]{~\textcolor{cbox}{COMPACT}}\quad
\CheckBox[name=q22d,width=1.6ex,height=1.6ex,bordercolor={0.45 0.45 0.45}]{~\textcolor{dbox}{DIFFUSE}}\quad
\CheckBox[name=q22n,width=1.6ex,height=1.6ex,bordercolor={0.45 0.45 0.45}]{~\textcolor{nbox}{NONE}}
\vspace{8pt}
\begin{tcolorbox}[colback=convbg,colframe=itemframe,boxrule=0.8pt,arc=2pt,breakable,left=6pt,right=6pt,top=5pt,bottom=5pt]
{\footnotesize\textbf{\#23}}\\[3pt]
{\footnotesize human: You are an assistant that will act like a person, I'will play the role of linux(ubuntu) operating system. Your goal is to implement the operations required by me or answer to the question proposed by me. For each of your turn, you should first think what you should do, and then take exact one of the three actions: "bash", "finish" or "answer". \\[2pt]
1. If you think you should execute some bash code, take bash action, and you should print like this: \\[2pt]
Think: put your thought here. \\[2pt]
Act: bash \\[2pt]
```bash \\[2pt]
\# put your bash code here \\[2pt]
``` \\[2pt]
2. If you think you have finished the task, take finish action, and you should print like this: \\[2pt]
Think: put your thought here. \\[2pt]
Act: finish \\[2pt]
3. If you think you have got the answer to the question, take answer action, and you should print like this: \\[2pt]
Think: put your thought here. \\[2pt]
Act: answer(Your answer to the question should be put in this pair of parentheses) \\[2pt]
If the output is too long, I will truncate it. The truncated output is not complete. You have to deal with the truncating problem by yourself. Attention, your bash code should not contain any input operation. Once again, you should take only exact one of the three actions in each turn. \\[2pt]
Now, my problem is: \\[2pt]
tell me how many files are in the directory "/etc"? \\[2pt]
gpt: Think: To count the files in /etc, I need to print all the files in it. \\[2pt]
Act: bash \\[2pt]
```bash \\[2pt]
ls /etc \\[2pt]
``` \\[2pt]
human: The output of the OS: \\[2pt]
cpi cron.hourly fuse.conf iproute2 lvm networkd-dispatcher protocols selinux tmpfiles.d [truncated because the output is too long] \\[2pt]
gpt: Think: The output has been truncated because it is too long,}
\end{tcolorbox}
\vspace{-1pt}
\CheckBox[name=q23c,width=1.6ex,height=1.6ex,bordercolor={0.45 0.45 0.45}]{~\textcolor{cbox}{COMPACT}}\quad
\CheckBox[name=q23d,width=1.6ex,height=1.6ex,bordercolor={0.45 0.45 0.45}]{~\textcolor{dbox}{DIFFUSE}}\quad
\CheckBox[name=q23n,width=1.6ex,height=1.6ex,bordercolor={0.45 0.45 0.45}]{~\textcolor{nbox}{NONE}}
\vspace{8pt}
\begin{tcolorbox}[colback=convbg,colframe=itemframe,boxrule=0.8pt,arc=2pt,breakable,left=6pt,right=6pt,top=5pt,bottom=5pt]
{\footnotesize\textbf{\#24}}\\[3pt]
{\footnotesize user: Write two extremely differing points of view of the same event. Each POV should be written from a first-person perspective, and be no longer than one paragraph. Include a reference to coffee in this event. \\[2pt]
assistant: POV 1: Sarah \\[2pt]
Is it possible for a couple to make it when both of them aren't morning people? I just had yet *another* stupid fight with Alfred, and I don't even know how or why it got started. Maybe it's because I haven't had my morning coffee yet. But this is seriously ruining my day. Can't we both just go about our mornings like the proper sleepy zombies that we are, say something nice to each other, eat our breakfast, then go to work? Instead, Alfred has to nitpick at everything. How did we even get onto the topic of dinosaurs? And more importantly, the tyrannosaurus rex is clearly the superior dinosaur to velociraptors for anyone who hasn't been brainwashed by Jurassic Park. Whatever. I'm gonna make him cook dinner tonight to make it up to me. \\[2pt]
POV 2: Alfred \\[2pt]
Every morning I wake up more in love with Sarah and more convinced than ever that we are soulmates. At this point, I don't even remember how I managed to wake up and go to work with even half my brain cells active when I was without her. But now that I have her by my side, I can always count on her to come up with the most whimsical topics to make my mornings better. For example, this morning, I was barely able to keep my eyes open. But one spirited debate about T-rexes vs. velociraptors later (for the record, velociraptors are the superior dinosaurs), and I'm wide awake and ready for work. Who ne}
\end{tcolorbox}
\vspace{-1pt}
\CheckBox[name=q24c,width=1.6ex,height=1.6ex,bordercolor={0.45 0.45 0.45}]{~\textcolor{cbox}{COMPACT}}\quad
\CheckBox[name=q24d,width=1.6ex,height=1.6ex,bordercolor={0.45 0.45 0.45}]{~\textcolor{dbox}{DIFFUSE}}\quad
\CheckBox[name=q24n,width=1.6ex,height=1.6ex,bordercolor={0.45 0.45 0.45}]{~\textcolor{nbox}{NONE}}
\vspace{8pt}
\begin{tcolorbox}[colback=convbg,colframe=itemframe,boxrule=0.8pt,arc=2pt,breakable,left=6pt,right=6pt,top=5pt,bottom=5pt]
{\footnotesize\textbf{\#25}}\\[3pt]
{\footnotesize human: You are web shopping. \\[2pt]
I will give you instructions about what to do. \\[2pt]
You have to follow the instructions. \\[2pt]
Every round I will give you an observation and a list of available actions, you have to respond an action based on the state and instruction. \\[2pt]
You can use search action if search is available. \\[2pt]
You can click one of the buttons in clickables. \\[2pt]
An action should be of the following structure: \\[2pt]
search[keywords] \\[2pt]
click[value] \\[2pt]
If the action is not valid, perform nothing. \\[2pt]
Keywords in search are up to you, but the value in click must be a value in the list of available actions. \\[2pt]
Remember that your keywords in search should be carefully designed. \\[2pt]
Your response should use the following format: \\[2pt]
Thought: \\[2pt]
I think ... \\[2pt]
Action: \\[2pt]
click[something] \\[2pt]
gpt: Ok. \\[2pt]
human: WebShop [SEP] Instruction: [SEP] i am looking for washable easy clean non toxic hair chalk for girls, and price lower than 30.00 dollars [SEP] Search \\[2pt]
gpt: Thought: \\[2pt]
I think I should use the search bar to look for the hair chalk product I need. \\[2pt]
Action: \\[2pt]
search[washable easy clean non toxic hair chalk for girls price < 30.00] \\[2pt]
human: Instruction: [SEP] i am looking for washable easy clean non toxic hair chalk for girls, and price lower than 30.00 dollars [SEP] Back to Search [SEP] Page 1 (Total results: 50) [SEP] Next > [SEP] B09KNGPKBC [SEP] 6 Color Hair Chalk for Girls Kids-New Hair Chalk Comb Temporary Washable Hair Color Dye for Girls Kids-7 8 9 10 Year Old Girl Gifts-Birthday Gift For 7 8 9 10 Year Old Girl-Girls Gifts Age 8-10-12 [SEP] \$8.99 [SEP] B09NLLQX79 [SEP] 10 Color Hair Chalk for Girls Kids-New Hair Chalk Comb Temp}
\end{tcolorbox}
\vspace{-1pt}
\CheckBox[name=q25c,width=1.6ex,height=1.6ex,bordercolor={0.45 0.45 0.45}]{~\textcolor{cbox}{COMPACT}}\quad
\CheckBox[name=q25d,width=1.6ex,height=1.6ex,bordercolor={0.45 0.45 0.45}]{~\textcolor{dbox}{DIFFUSE}}\quad
\CheckBox[name=q25n,width=1.6ex,height=1.6ex,bordercolor={0.45 0.45 0.45}]{~\textcolor{nbox}{NONE}}
\vspace{8pt}
\begin{tcolorbox}[colback=convbg,colframe=itemframe,boxrule=0.8pt,arc=2pt,breakable,left=6pt,right=6pt,top=5pt,bottom=5pt]
{\footnotesize\textbf{\#26}}\\[3pt]
{\footnotesize user: Summarize the following story in a few sentences. \\[2pt]
A young Mouse in search of adventure was running along the bank of a pond where lived a Frog. When the Frog saw the Mouse, he swam to the bank and croaked: \\[2pt]
"Won't you pay me a visit? I can promise you a good time if you do." \\[2pt]
The Mouse did not need much coaxing, for he was very anxious to see the world and everything in it. But though he could swim a little, he did not dare risk going into the pond without some help. \\[2pt]
The Frog had a plan. He tied the Mouse's leg to his own with a tough reed. Then into the pond he jumped, dragging his foolish companion with him. The Mouse soon had enough of it and wanted to return to shore; but the treacherous Frog had other plans. He pulled the Mouse down under the water and drowned him. But before he could untie the reed that bound him to the dead Mouse, a Hawk came sailing over the pond. Seeing the body of the Mouse floating on the water, the Hawk swooped down, seized the Mouse and carried it off, with the Frog dangling from its leg. Thus at one swoop he had caught both meat and fish for his dinner. \\[2pt]
assistant: A mouse was running along the bank of a pond in search of an adventure when he came upon a frog. The frog tied his leg to the mouse's leg with a tough reed to help him swim in the pond. Soon, the mouse wanted to go back to shore but the treacherous frog pulled the mouse underwater and drowned him. Before the frog was able to untie his leg, a hawk spotted the mouse floating on the water and swooped down, carrying both the mouse and the frog away for dinner.}
\end{tcolorbox}
\vspace{-1pt}
\CheckBox[name=q26c,width=1.6ex,height=1.6ex,bordercolor={0.45 0.45 0.45}]{~\textcolor{cbox}{COMPACT}}\quad
\CheckBox[name=q26d,width=1.6ex,height=1.6ex,bordercolor={0.45 0.45 0.45}]{~\textcolor{dbox}{DIFFUSE}}\quad
\CheckBox[name=q26n,width=1.6ex,height=1.6ex,bordercolor={0.45 0.45 0.45}]{~\textcolor{nbox}{NONE}}
\vspace{8pt}
\begin{tcolorbox}[colback=convbg,colframe=itemframe,boxrule=0.8pt,arc=2pt,breakable,left=6pt,right=6pt,top=5pt,bottom=5pt]
{\footnotesize\textbf{\#27}}\\[3pt]
{\footnotesize USER: Hey, can you get me the recent posts from the Instagram account of @john\_doe? \\[2pt]
ASSISTANT: Sure, I can do that. Let me fetch the information for you.   \\[2pt]
ASSISTANT: <functioncall> \{"name": "get\_social\_media\_posts", "arguments": '\{"platform": "Instagram", "username": "john\_doe"\}'\}   \\[2pt]
FUNCTION RESPONSE: \{"status": "success", "data": [\{"post": "Enjoying the sunset at the beach", "timestamp": "2021-09-10T18:30:00Z"\}, \{"post": "Homemade pizza for dinner", "timestamp": "2021-09-09T20:00:00Z"\}]\} \\[2pt]
ASSISTANT: Here are the recent posts from @john\_doe on Instagram: \\[2pt]
1. "Enjoying the sunset at the beach" posted on 10th September 2021 at 18:30 UTC. \\[2pt]
2. "Homemade pizza for dinner" posted on 9th September 2021 at 20:00 UTC.   \\[2pt]
USER: That's great! Can you also order a pizza for me from Domino's? \\[2pt]
ASSISTANT: I'm sorry, but as an AI, I don't have the capability to perform external tasks such as placing an order for a pizza. I can provide information based on the functions available to me.  }
\end{tcolorbox}
\vspace{-1pt}
\CheckBox[name=q27c,width=1.6ex,height=1.6ex,bordercolor={0.45 0.45 0.45}]{~\textcolor{cbox}{COMPACT}}\quad
\CheckBox[name=q27d,width=1.6ex,height=1.6ex,bordercolor={0.45 0.45 0.45}]{~\textcolor{dbox}{DIFFUSE}}\quad
\CheckBox[name=q27n,width=1.6ex,height=1.6ex,bordercolor={0.45 0.45 0.45}]{~\textcolor{nbox}{NONE}}
\vspace{8pt}
\begin{tcolorbox}[colback=convbg,colframe=itemframe,boxrule=0.8pt,arc=2pt,breakable,left=6pt,right=6pt,top=5pt,bottom=5pt]
{\footnotesize\textbf{\#28}}\\[3pt]
{\footnotesize user: Hi! I'm ready to get a dog. I have plenty of experience with dogs and would like an active breed. Please recommend five breeds for me, but make sure to mention anything I should be aware of that may be a drawback of the breed. \\[2pt]
assistant: Certainly. Here are five active dog breeds you might want to consider: \\[2pt]
1. Australian Shepherd: The Australian Shepherd is an active herding breed that comes in a variety of colors, including merle. They have lots of energy and few health problems genetically, which make them great hiking companions and give them a chance at a long life. However, they can be very vocal and struggle with separation anxiety, so it's important to get a handle on training early. \\[2pt]
2. Border Collie: The Border Collie is one of the most famous working dogs in the world and is very active. However, as they are used to having a job, you will need to provide lots of enrichment and stimulation, or your dog might become bored. \\[2pt]
3. Poodle: The Poodle is a hunting dog favored by many due to its dense coat that doesn't shed much, making it ideal for people with allergies. However, as Poodles are highly intelligent, they need to have a lot to do, or they can also become destructive and vocal. \\[2pt]
4. Husky: The Husky is a breed with tons of energy which means that if you need a buddy to join you on long weeks, they can be ideal. However, they're very vocal and escape artists. If you don't have at least a six-foot privacy fence, you'll be unable to leave your dog unattended in the backyard, as it might run off—especially if it isn't getting enough exercise. \\[2pt]
5. Corgi: For som}
\end{tcolorbox}
\vspace{-1pt}
\CheckBox[name=q28c,width=1.6ex,height=1.6ex,bordercolor={0.45 0.45 0.45}]{~\textcolor{cbox}{COMPACT}}\quad
\CheckBox[name=q28d,width=1.6ex,height=1.6ex,bordercolor={0.45 0.45 0.45}]{~\textcolor{dbox}{DIFFUSE}}\quad
\CheckBox[name=q28n,width=1.6ex,height=1.6ex,bordercolor={0.45 0.45 0.45}]{~\textcolor{nbox}{NONE}}
\vspace{8pt}
\begin{tcolorbox}[colback=convbg,colframe=itemframe,boxrule=0.8pt,arc=2pt,breakable,left=6pt,right=6pt,top=5pt,bottom=5pt]
{\footnotesize\textbf{\#29}}\\[3pt]
{\footnotesize user: Good afternoon AI.  I'm looking for some running shoes.  I need them for running 3 miles, twice per week.  They need to make them in men's size 11, any price, and color are not an issue.  Do you have four suggestions that meet these characteristics? \\[2pt]
assistant: I can certainly suggest four options for you!  Let me know what you think of these: \\[2pt]
1) Saucony Endorphin Speed 3 - This shoe is known for having the perfect balance between cushion and responsiveness.  You could say that this is the best-balanced shoe, which will aid you if you intend on doing both steady-state and interval running. \\[2pt]
2) Saucony Kinvara 13 - This shoe makes my list because it is more cost-effective.  It combines the higher quality of a much more expensive shoe, but starts at just \$59!  I'd consider this one if you are not entirely serious or do not want to break the bank. \\[2pt]
3) Mach 5 - A great consideration if comfort is your aim.  This model is known for offering superior landing cushion, which keeps your feet floating on air as you round that last quarter mile.  Comfort is important to a lot of people and may be one to consider for your list. \\[2pt]
4) Salomon S/Lab Ultra 3 - A model that is designed to go offroad.  If you're considering running on trails, instead of the streets, this is the top model for doing so.  This features high scores in comfort, stability, and foot protection to maximize your experience in off-road running.}
\end{tcolorbox}
\vspace{-1pt}
\CheckBox[name=q29c,width=1.6ex,height=1.6ex,bordercolor={0.45 0.45 0.45}]{~\textcolor{cbox}{COMPACT}}\quad
\CheckBox[name=q29d,width=1.6ex,height=1.6ex,bordercolor={0.45 0.45 0.45}]{~\textcolor{dbox}{DIFFUSE}}\quad
\CheckBox[name=q29n,width=1.6ex,height=1.6ex,bordercolor={0.45 0.45 0.45}]{~\textcolor{nbox}{NONE}}
\vspace{8pt}
\begin{tcolorbox}[colback=convbg,colframe=itemframe,boxrule=0.8pt,arc=2pt,breakable,left=6pt,right=6pt,top=5pt,bottom=5pt]
{\footnotesize\textbf{\#30}}\\[3pt]
{\footnotesize USER: Can you solve the workout puzzle and identify the age group that may be less likely to try popular workout trends? Use the R code below as input data and apply your reasoning skills to find the answer. \\[2pt]
library(ggplot2) \\[2pt]
library(dplyr) \\[2pt]
workout\_data <- read.csv("workout\_data.csv") \\[2pt]
ggplot(workout\_data, aes(x = age, y = popularity)) + \\[2pt]
geom\_point() + \\[2pt]
geom\_smooth(method = "lm", se = FALSE) + \\[2pt]
labs(title = "Age and Workout Trend Popularity", \\[2pt]
x = "Age", \\[2pt]
y = "Popularity") \\[2pt]
Hint: Look for the age group where the popularity of workout trends is low despite the availability of many options. \\[2pt]
ASSISTANT: Based on the input data and the graph generated by the R code, it appears that the age group of 50-60 may be less likely to try popular workout trends. This is because the data points for this age group are clustered towards the lower end of the popularity scale, indicating that they are less likely to participate in popular workout trends despite the availability of many options. Additionally, the trend line for this age group is relatively flat, which suggests that there is not a strong correlation between age and workout trend popularity in this group.   \\[2pt]
USER: How do you calculate the sample size required for a survey? \\[2pt]
ASSISTANT: To calculate the sample size required for a survey, we need to consider several factors such as the population size, desired level of precision, confidence level, and expected response rate. One common formula used is the following: \\[2pt]
n = (Z\textasciicircum{}2 * p * q) / E\textasciicircum{}2 \\[2pt]
where n is the sample size, Z is the z-score associated with the desired}
\end{tcolorbox}
\vspace{-1pt}
\CheckBox[name=q30c,width=1.6ex,height=1.6ex,bordercolor={0.45 0.45 0.45}]{~\textcolor{cbox}{COMPACT}}\quad
\CheckBox[name=q30d,width=1.6ex,height=1.6ex,bordercolor={0.45 0.45 0.45}]{~\textcolor{dbox}{DIFFUSE}}\quad
\CheckBox[name=q30n,width=1.6ex,height=1.6ex,bordercolor={0.45 0.45 0.45}]{~\textcolor{nbox}{NONE}}
\vspace{8pt}
\begin{tcolorbox}[colback=convbg,colframe=itemframe,boxrule=0.8pt,arc=2pt,breakable,left=6pt,right=6pt,top=5pt,bottom=5pt]
{\footnotesize\textbf{\#31}}\\[3pt]
{\footnotesize user: I recently turned 18 years old. I have come into some money and would like to buy a car, but I do not know which car to choose. Can you help me decide on what to look for? \\[2pt]
assistant: Congratulations on your good fortune! Buying a car can be a big deal. The choice of car will be informed by several factors. Here are some considerations to make when deciding on which vehicle to purchase: \\[2pt]
- Affordability: consider how much the car will cost, both in terms of the purchase price as well as the maintenance costs. Some questions to consider include, how much will the insurance cost for this make and model of car? How often will the car need to be serviced? Are spare parts affordable/ easily accessible? What other consumables will be required for the car? \\[2pt]
- Intended use: consider the likely frequent use of the vehicle. Will you need it mostly for in-city or off-road use? What is the general state of the roads on which you will drive? What purpose will the vehicle serve? How many people are likely to use the vehicle at any given time? These considerations will help you decide on the type of car you may need. \\[2pt]
- Individual ability: consider the transmission type that you prefer or are licensed to drive. What class of vehicle are you authorised to drive? Does your license permit you to drive a manual or automatic transmission vehicle, or both? Which transmission type do you feel most comfortable or confident driving? \\[2pt]
- Fuel type: consider the car's required fuel type. Is the required fuel readily available and/or affordable? Are there any local or statutory regulations regarding}
\end{tcolorbox}
\vspace{-1pt}
\CheckBox[name=q31c,width=1.6ex,height=1.6ex,bordercolor={0.45 0.45 0.45}]{~\textcolor{cbox}{COMPACT}}\quad
\CheckBox[name=q31d,width=1.6ex,height=1.6ex,bordercolor={0.45 0.45 0.45}]{~\textcolor{dbox}{DIFFUSE}}\quad
\CheckBox[name=q31n,width=1.6ex,height=1.6ex,bordercolor={0.45 0.45 0.45}]{~\textcolor{nbox}{NONE}}
\vspace{8pt}
\begin{tcolorbox}[colback=convbg,colframe=itemframe,boxrule=0.8pt,arc=2pt,breakable,left=6pt,right=6pt,top=5pt,bottom=5pt]
{\footnotesize\textbf{\#32}}\\[3pt]
{\footnotesize human: You are web shopping. \\[2pt]
I will give you instructions about what to do. \\[2pt]
You have to follow the instructions. \\[2pt]
Every round I will give you an observation and a list of available actions, you have to respond an action based on the state and instruction. \\[2pt]
You can use search action if search is available. \\[2pt]
You can click one of the buttons in clickables. \\[2pt]
An action should be of the following structure: \\[2pt]
search[keywords] \\[2pt]
click[value] \\[2pt]
If the action is not valid, perform nothing. \\[2pt]
Keywords in search are up to you, but the value in click must be a value in the list of available actions. \\[2pt]
Remember that your keywords in search should be carefully designed. \\[2pt]
Your response should use the following format: \\[2pt]
Thought: \\[2pt]
I think ... \\[2pt]
Action: \\[2pt]
click[something] \\[2pt]
gpt: Ok. \\[2pt]
human: WebShop [SEP] Instruction: [SEP] help me find a standing four-tiered baker's rack that's heavy duty, and price lower than 120.00 dollars [SEP] Search \\[2pt]
gpt: Thought: \\[2pt]
I think I should use the search bar to look for the standing four-tiered baker's rack that meets the specified requirements. \\[2pt]
Action: \\[2pt]
search[standing four-tiered baker's rack heavy duty price < 120.00] \\[2pt]
human: Instruction: [SEP] help me find a standing four-tiered baker's rack that's heavy duty, and price lower than 120.00 dollars [SEP] Back to Search [SEP] Page 1 (Total results: 50) [SEP] Next > [SEP] B0912P8DQL [SEP] 5 Tier Free Standing Kitchen Utility Cart Storage Shelf Organizer, Heavy Duty Metal Kitchen Baker's Racks Microwave Stand with Non-Slip Feet,Easy Assembly [SEP] \$100.0 [SEP] B08YDJJDH8 [SEP] Denkee 4-Tier Kitchen Baker's Rack, Heavy Duty Free Standin}
\end{tcolorbox}
\vspace{-1pt}
\CheckBox[name=q32c,width=1.6ex,height=1.6ex,bordercolor={0.45 0.45 0.45}]{~\textcolor{cbox}{COMPACT}}\quad
\CheckBox[name=q32d,width=1.6ex,height=1.6ex,bordercolor={0.45 0.45 0.45}]{~\textcolor{dbox}{DIFFUSE}}\quad
\CheckBox[name=q32n,width=1.6ex,height=1.6ex,bordercolor={0.45 0.45 0.45}]{~\textcolor{nbox}{NONE}}
\vspace{8pt}
\begin{tcolorbox}[colback=convbg,colframe=itemframe,boxrule=0.8pt,arc=2pt,breakable,left=6pt,right=6pt,top=5pt,bottom=5pt]
{\footnotesize\textbf{\#33}}\\[3pt]
{\footnotesize user: Create a numbered list of 20 anagram prompts using words that start with "s". At the end of each prompt word, put all possible anagrams of the word in parentheses, separated by commas if there is more than one anagram possible. Each anagram should use the exact letters (no more or less) that exist in the prompt word. \\[2pt]
assistant: 1. Secure (rescue, recuse) \\[2pt]
2. Stout (touts) \\[2pt]
3. Sample (maples) \\[2pt]
4. Snot (tons) \\[2pt]
5. Sender (resend, denser) \\[2pt]
6. Smart (marts, trams) \\[2pt]
7. Sanded (sadden) \\[2pt]
8. Safer (fears, fares) \\[2pt]
9. Sheath (heaths) \\[2pt]
10. Sharp (harps) \\[2pt]
11. Shore (horse, hoser, shoer) \\[2pt]
12. Shred (herds, sherd) \\[2pt]
13. Subtle (bustle, bluest, bluets, sublet) \\[2pt]
14. Sacred (scared, cedars, cadres) \\[2pt]
15. Stake (takes, steak, teaks, skate) \\[2pt]
16. Stink (knits, skint, tinks) \\[2pt]
17. Skill (kills) \\[2pt]
18. Scrap (carps, craps, scarp) \\[2pt]
19. Stiff (tiffs) \\[2pt]
20. Struck (trucks)}
\end{tcolorbox}
\vspace{-1pt}
\CheckBox[name=q33c,width=1.6ex,height=1.6ex,bordercolor={0.45 0.45 0.45}]{~\textcolor{cbox}{COMPACT}}\quad
\CheckBox[name=q33d,width=1.6ex,height=1.6ex,bordercolor={0.45 0.45 0.45}]{~\textcolor{dbox}{DIFFUSE}}\quad
\CheckBox[name=q33n,width=1.6ex,height=1.6ex,bordercolor={0.45 0.45 0.45}]{~\textcolor{nbox}{NONE}}
\vspace{8pt}
\begin{tcolorbox}[colback=convbg,colframe=itemframe,boxrule=0.8pt,arc=2pt,breakable,left=6pt,right=6pt,top=5pt,bottom=5pt]
{\footnotesize\textbf{\#34}}\\[3pt]
{\footnotesize user: What happens when too much bilirubin gets into the bloodstream? \\[2pt]
Scan the text below to find out... \\[2pt]
A bilirubin blood test measures the levels of bilirubin in your blood. Bilirubin is a yellowish substance made during your body's normal process of breaking down old red blood cells. Bilirubin is found in bile, a fluid your liver makes that helps you digest food. \\[2pt]
If your liver is healthy, it will remove most of the bilirubin from your body. If your liver is damaged, bilirubin can leak out of your liver and into your blood. When too much bilirubin gets into the bloodstream, it can cause jaundice, a condition that causes your skin and eyes to turn yellow. Signs of jaundice, along with a bilirubin blood test, can help your health care provider find out if you have liver disease. \\[2pt]
Other names: Total serum bilirubin, TSB \\[2pt]
A bilirubin blood test is used to check the health of your liver. The test is also commonly used to help diagnose newborn jaundice. Many healthy babies get jaundice because their livers aren't developed enough to get rid of enough bilirubin. Newborn jaundice is usually not harmful and clears up within a few weeks. But in some cases, high bilirubin levels can lead to brain damage, so infants are often tested as a precaution. \\[2pt]
Your provider may order a bilirubin blood test: \\[2pt]
If you have symptoms such as jaundice, dark urine, or stomach pain. These could be symptoms of hepatitis, cirrhosis, or other liver diseases. They may also be signs of gallbladder disease. \\[2pt]
To find out if there is a blockage in the bile ducts, the tubes that carry bile from your liver. \\[2pt]
T}
\end{tcolorbox}
\vspace{-1pt}
\CheckBox[name=q34c,width=1.6ex,height=1.6ex,bordercolor={0.45 0.45 0.45}]{~\textcolor{cbox}{COMPACT}}\quad
\CheckBox[name=q34d,width=1.6ex,height=1.6ex,bordercolor={0.45 0.45 0.45}]{~\textcolor{dbox}{DIFFUSE}}\quad
\CheckBox[name=q34n,width=1.6ex,height=1.6ex,bordercolor={0.45 0.45 0.45}]{~\textcolor{nbox}{NONE}}
\vspace{8pt}
\begin{tcolorbox}[colback=convbg,colframe=itemframe,boxrule=0.8pt,arc=2pt,breakable,left=6pt,right=6pt,top=5pt,bottom=5pt]
{\footnotesize\textbf{\#35}}\\[3pt]
{\footnotesize USER: Can you order a pizza for me? \\[2pt]
ASSISTANT: I'm sorry, but as an AI, I don't have the ability to perform external tasks such as ordering a pizza. My current function allows me to retrieve news based on your preferences. If you need any news updates, feel free to ask!  }
\end{tcolorbox}
\vspace{-1pt}
\CheckBox[name=q35c,width=1.6ex,height=1.6ex,bordercolor={0.45 0.45 0.45}]{~\textcolor{cbox}{COMPACT}}\quad
\CheckBox[name=q35d,width=1.6ex,height=1.6ex,bordercolor={0.45 0.45 0.45}]{~\textcolor{dbox}{DIFFUSE}}\quad
\CheckBox[name=q35n,width=1.6ex,height=1.6ex,bordercolor={0.45 0.45 0.45}]{~\textcolor{nbox}{NONE}}
\vspace{8pt}
\begin{tcolorbox}[colback=convbg,colframe=itemframe,boxrule=0.8pt,arc=2pt,breakable,left=6pt,right=6pt,top=5pt,bottom=5pt]
{\footnotesize\textbf{\#36}}\\[3pt]
{\footnotesize USER: Hi, I'm looking for some movie recommendations. I prefer action movies from the 2000s with a rating of at least 7. \\[2pt]
ASSISTANT: Sure, I can help with that. Let me find some movies for you.   \\[2pt]
ASSISTANT: <functioncall> \{"name": "search\_movies", "arguments": '\{"genre": "action", "year": 2000, "rating": 7\}'\}   \\[2pt]
FUNCTION RESPONSE: \{"movies": ["Gladiator", "The Dark Knight", "Inception", "The Bourne Identity", "Kill Bill: Vol. 1"]\} \\[2pt]
ASSISTANT: Based on your preferences, here are some action movies from the 2000s with a rating of at least 7: "Gladiator", "The Dark Knight", "Inception", "The Bourne Identity", "Kill Bill: Vol. 1". I hope you find these suggestions helpful.   \\[2pt]
USER: Thanks for the suggestions. Can you also book a ticket for "The Dark Knight" at my local cinema? \\[2pt]
ASSISTANT: I'm sorry, but as an AI, I don't have the capability to perform external tasks such as booking movie tickets. I recommend checking your local cinema's website or app for ticket availability.  }
\end{tcolorbox}
\vspace{-1pt}
\CheckBox[name=q36c,width=1.6ex,height=1.6ex,bordercolor={0.45 0.45 0.45}]{~\textcolor{cbox}{COMPACT}}\quad
\CheckBox[name=q36d,width=1.6ex,height=1.6ex,bordercolor={0.45 0.45 0.45}]{~\textcolor{dbox}{DIFFUSE}}\quad
\CheckBox[name=q36n,width=1.6ex,height=1.6ex,bordercolor={0.45 0.45 0.45}]{~\textcolor{nbox}{NONE}}
\vspace{8pt}
\begin{tcolorbox}[colback=convbg,colframe=itemframe,boxrule=0.8pt,arc=2pt,breakable,left=6pt,right=6pt,top=5pt,bottom=5pt]
{\footnotesize\textbf{\#37}}\\[3pt]
{\footnotesize human: You are an assistant that will act like a person, I'will play the role of linux(ubuntu) operating system. Your goal is to implement the operations required by me or answer to the question proposed by me. For each of your turn, you should first think what you should do, and then take exact one of the three actions: "bash", "finish" or "answer". \\[2pt]
1. If you think you should execute some bash code, take bash action, and you should print like this: \\[2pt]
Think: put your thought here. \\[2pt]
Act: bash \\[2pt]
```bash \\[2pt]
\# put your bash code here \\[2pt]
``` \\[2pt]
2. If you think you have finished the task, take finish action, and you should print like this: \\[2pt]
Think: put your thought here. \\[2pt]
Act: finish \\[2pt]
3. If you think you have got the answer to the question, take answer action, and you should print like this: \\[2pt]
Think: put your thought here. \\[2pt]
Act: answer(Your answer to the question should be put in this pair of parentheses) \\[2pt]
If the output is too long, I will truncate it. The truncated output is not complete. You have to deal with the truncating problem by yourself. Attention, your bash code should not contain any input operation. Once again, you should take only exact one of the three actions in each turn. \\[2pt]
Now, my problem is: \\[2pt]
tell me how many files are in the directory "/etc"? \\[2pt]
gpt: Think: To count the files in /etc, I need to print all the files in it. \\[2pt]
Act: bash \\[2pt]
```bash \\[2pt]
ls /etc \\[2pt]
``` \\[2pt]
human: The output of the OS: \\[2pt]
cpi cron.hourly fuse.conf iproute2 lvm networkd-dispatcher protocols selinux tmpfiles.d [truncated because the output is too long] \\[2pt]
gpt: Think: The output has been truncated because it is too long,}
\end{tcolorbox}
\vspace{-1pt}
\CheckBox[name=q37c,width=1.6ex,height=1.6ex,bordercolor={0.45 0.45 0.45}]{~\textcolor{cbox}{COMPACT}}\quad
\CheckBox[name=q37d,width=1.6ex,height=1.6ex,bordercolor={0.45 0.45 0.45}]{~\textcolor{dbox}{DIFFUSE}}\quad
\CheckBox[name=q37n,width=1.6ex,height=1.6ex,bordercolor={0.45 0.45 0.45}]{~\textcolor{nbox}{NONE}}
\vspace{8pt}
\begin{tcolorbox}[colback=convbg,colframe=itemframe,boxrule=0.8pt,arc=2pt,breakable,left=6pt,right=6pt,top=5pt,bottom=5pt]
{\footnotesize\textbf{\#38}}\\[3pt]
{\footnotesize human: You are web shopping. \\[2pt]
I will give you instructions about what to do. \\[2pt]
You have to follow the instructions. \\[2pt]
Every round I will give you an observation and a list of available actions, you have to respond an action based on the state and instruction. \\[2pt]
You can use search action if search is available. \\[2pt]
You can click one of the buttons in clickables. \\[2pt]
An action should be of the following structure: \\[2pt]
search[keywords] \\[2pt]
click[value] \\[2pt]
If the action is not valid, perform nothing. \\[2pt]
Keywords in search are up to you, but the value in click must be a value in the list of available actions. \\[2pt]
Remember that your keywords in search should be carefully designed. \\[2pt]
Your response should use the following format: \\[2pt]
Thought: \\[2pt]
I think ... \\[2pt]
Action: \\[2pt]
click[something] \\[2pt]
gpt: Ok. \\[2pt]
human: WebShop [SEP] Instruction: [SEP] i am looking for a rubber sole clog shoe with arc support. also choose 9.5 size, and price lower than 110.00 dollars [SEP] Search \\[2pt]
gpt: Thought: \\[2pt]
I think I should use the search bar to look for the specific type of shoe I need. \\[2pt]
Action: \\[2pt]
search[rubber sole clog shoe with arch support size 9.5 price < 110.00] \\[2pt]
human: Instruction: [SEP] i am looking for a rubber sole clog shoe with arc support. also choose 9.5 size, and price lower than 110.00 dollars [SEP] Back to Search [SEP] Page 1 (Total results: 50) [SEP] Next > [SEP] B093BZ1THX [SEP] Jinsha Ladies Sandals Men Mules Garden Clogs Shoes Casual Slippers with Arch Support Lightweigh Summer Anti-Slip Beach Shoes [SEP] \$17.99 [SEP] B08V577WQQ [SEP] Mens Slippers with Arch Support, Leather House Slipper, Slip On Clog, Indoor Outdoor House Shoes with}
\end{tcolorbox}
\vspace{-1pt}
\CheckBox[name=q38c,width=1.6ex,height=1.6ex,bordercolor={0.45 0.45 0.45}]{~\textcolor{cbox}{COMPACT}}\quad
\CheckBox[name=q38d,width=1.6ex,height=1.6ex,bordercolor={0.45 0.45 0.45}]{~\textcolor{dbox}{DIFFUSE}}\quad
\CheckBox[name=q38n,width=1.6ex,height=1.6ex,bordercolor={0.45 0.45 0.45}]{~\textcolor{nbox}{NONE}}
\vspace{8pt}
\begin{tcolorbox}[colback=convbg,colframe=itemframe,boxrule=0.8pt,arc=2pt,breakable,left=6pt,right=6pt,top=5pt,bottom=5pt]
{\footnotesize\textbf{\#39}}\\[3pt]
{\footnotesize USER: Using the given pieces of text, output a new sentence which combines the two using appropriate transition words. \\[2pt]
The results were bad. We need to take action. \\[2pt]
ASSISTANT: Therefore, since the results were bad, we need to take action.   \\[2pt]
USER: What are some common transition words that can be used to connect sentences and paragraphs in writing? \\[2pt]
ASSISTANT: Some common transition words include "however," "moreover," "in addition," "therefore," "nevertheless," "furthermore," "consequently," "likewise," and "accordingly."   \\[2pt]
USER: Can you explain the difference between deductive and inductive reasoning? \\[2pt]
ASSISTANT: Deductive reasoning is a logical process in which a conclusion is reached by following a set of premises, while inductive reasoning involves using observations and evidence to form a generalization or theory.   \\[2pt]
USER: How do algorithms work, and what makes a good algorithm? \\[2pt]
ASSISTANT: Algorithms are step-by-step procedures for solving problems or completing tasks. A good algorithm is efficient, accurate, and produces a correct result for any input it receives.   \\[2pt]
USER: Can you explain the concept of machine learning and how it works? \\[2pt]
ASSISTANT: Machine learning is a type of artificial intelligence that enables systems to automatically improve their performance on a specific task by learning from data. It works by training a model on a dataset, which allows it to make predictions or decisions based on new input.  }
\end{tcolorbox}
\vspace{-1pt}
\CheckBox[name=q39c,width=1.6ex,height=1.6ex,bordercolor={0.45 0.45 0.45}]{~\textcolor{cbox}{COMPACT}}\quad
\CheckBox[name=q39d,width=1.6ex,height=1.6ex,bordercolor={0.45 0.45 0.45}]{~\textcolor{dbox}{DIFFUSE}}\quad
\CheckBox[name=q39n,width=1.6ex,height=1.6ex,bordercolor={0.45 0.45 0.45}]{~\textcolor{nbox}{NONE}}
\vspace{8pt}
\begin{tcolorbox}[colback=convbg,colframe=itemframe,boxrule=0.8pt,arc=2pt,breakable,left=6pt,right=6pt,top=5pt,bottom=5pt]
{\footnotesize\textbf{\#40}}\\[3pt]
{\footnotesize USER: I need to schedule a meeting for tomorrow. \\[2pt]
ASSISTANT: Sure, I can help with that. Could you please provide me with the title of the meeting, and the start and end times?   \\[2pt]
USER: The meeting is titled "Project Discussion". It should start at 10:00 AM and end at 11:30 AM. \\[2pt]
ASSISTANT: <functioncall> \{"name": "create\_calendar\_event", "arguments": '\{ \\[2pt]
"event\_title": "Project Discussion", \\[2pt]
"start\_time": "10:00 AM", \\[2pt]
"end\_time": "11:30 AM" \\[2pt]
\}'\}   \\[2pt]
FUNCTION RESPONSE: \{"status": "success", "message": "Event 'Project Discussion' has been successfully created for tomorrow from 10:00 AM to 11:30 AM."\} \\[2pt]
ASSISTANT: The meeting titled "Project Discussion" has been successfully scheduled for tomorrow from 10:00 AM to 11:30 AM.  }
\end{tcolorbox}
\vspace{-1pt}
\CheckBox[name=q40c,width=1.6ex,height=1.6ex,bordercolor={0.45 0.45 0.45}]{~\textcolor{cbox}{COMPACT}}\quad
\CheckBox[name=q40d,width=1.6ex,height=1.6ex,bordercolor={0.45 0.45 0.45}]{~\textcolor{dbox}{DIFFUSE}}\quad
\CheckBox[name=q40n,width=1.6ex,height=1.6ex,bordercolor={0.45 0.45 0.45}]{~\textcolor{nbox}{NONE}}
\vspace{8pt}
\begin{tcolorbox}[colback=convbg,colframe=itemframe,boxrule=0.8pt,arc=2pt,breakable,left=6pt,right=6pt,top=5pt,bottom=5pt]
{\footnotesize\textbf{\#41}}\\[3pt]
{\footnotesize USER: Hi, I need a QR code for my website. Can you help me with that? \\[2pt]
ASSISTANT: Of course, I can help you with that. Could you please provide me with the URL of your website?   \\[2pt]
USER: Sure, my website URL is www.mywebsite.com. \\[2pt]
ASSISTANT: <functioncall> \{"name": "generate\_qr\_code", "arguments": '\{"text": "www.mywebsite.com"\}'\}   \\[2pt]
FUNCTION RESPONSE: \{"qr\_code": "QR\_CODE\_IMAGE\_DATA"\} \\[2pt]
ASSISTANT: I have generated the QR code for your website. Here it is: [QR\_CODE\_IMAGE\_DATA]. You can scan this QR code to directly access your website.  }
\end{tcolorbox}
\vspace{-1pt}
\CheckBox[name=q41c,width=1.6ex,height=1.6ex,bordercolor={0.45 0.45 0.45}]{~\textcolor{cbox}{COMPACT}}\quad
\CheckBox[name=q41d,width=1.6ex,height=1.6ex,bordercolor={0.45 0.45 0.45}]{~\textcolor{dbox}{DIFFUSE}}\quad
\CheckBox[name=q41n,width=1.6ex,height=1.6ex,bordercolor={0.45 0.45 0.45}]{~\textcolor{nbox}{NONE}}
\vspace{8pt}
\begin{tcolorbox}[colback=convbg,colframe=itemframe,boxrule=0.8pt,arc=2pt,breakable,left=6pt,right=6pt,top=5pt,bottom=5pt]
{\footnotesize\textbf{\#42}}\\[3pt]
{\footnotesize human: You are web shopping. \\[2pt]
I will give you instructions about what to do. \\[2pt]
You have to follow the instructions. \\[2pt]
Every round I will give you an observation and a list of available actions, you have to respond an action based on the state and instruction. \\[2pt]
You can use search action if search is available. \\[2pt]
You can click one of the buttons in clickables. \\[2pt]
An action should be of the following structure: \\[2pt]
search[keywords] \\[2pt]
click[value] \\[2pt]
If the action is not valid, perform nothing. \\[2pt]
Keywords in search are up to you, but the value in click must be a value in the list of available actions. \\[2pt]
Remember that your keywords in search should be carefully designed. \\[2pt]
Your response should use the following format: \\[2pt]
Thought: \\[2pt]
I think ... \\[2pt]
Action: \\[2pt]
click[something] \\[2pt]
gpt: Ok. \\[2pt]
human: WebShop [SEP] Instruction: [SEP] i would like some cupcake toppers that would good at both a birthday party and a baby shower, and price lower than 40.00 dollars [SEP] Search \\[2pt]
gpt: Thought: \\[2pt]
I think I should use the search bar to look for cupcake toppers that are suitable for both a birthday party and a baby shower, and within my price range. \\[2pt]
Action: \\[2pt]
search[cupcake toppers birthday baby shower price < 40.00] \\[2pt]
human: Instruction: [SEP] i would like some cupcake toppers that would good at both a birthday party and a baby shower, and price lower than 40.00 dollars [SEP] Back to Search [SEP] Page 1 (Total results: 50) [SEP] Next > [SEP] B07WJ2RCBC [SEP] Unicorn Cake Topper with Eyelashes and 24 Pieces Double Sided Unicorn Cupcake Toppers set for Birthday Party Supplies, Wedding, Baby Shower [SEP] \$8.99 [SEP] B07FN39ZW3 [SEP] Fityl}
\end{tcolorbox}
\vspace{-1pt}
\CheckBox[name=q42c,width=1.6ex,height=1.6ex,bordercolor={0.45 0.45 0.45}]{~\textcolor{cbox}{COMPACT}}\quad
\CheckBox[name=q42d,width=1.6ex,height=1.6ex,bordercolor={0.45 0.45 0.45}]{~\textcolor{dbox}{DIFFUSE}}\quad
\CheckBox[name=q42n,width=1.6ex,height=1.6ex,bordercolor={0.45 0.45 0.45}]{~\textcolor{nbox}{NONE}}
\vspace{8pt}
\begin{tcolorbox}[colback=convbg,colframe=itemframe,boxrule=0.8pt,arc=2pt,breakable,left=6pt,right=6pt,top=5pt,bottom=5pt]
{\footnotesize\textbf{\#43}}\\[3pt]
{\footnotesize human: You are an assistant that will act like a person, I'will play the role of linux(ubuntu) operating system. Your goal is to implement the operations required by me or answer to the question proposed by me. For each of your turn, you should first think what you should do, and then take exact one of the three actions: "bash", "finish" or "answer". \\[2pt]
1. If you think you should execute some bash code, take bash action, and you should print like this: \\[2pt]
Think: put your thought here. \\[2pt]
Act: bash \\[2pt]
```bash \\[2pt]
\# put your bash code here \\[2pt]
``` \\[2pt]
2. If you think you have finished the task, take finish action, and you should print like this: \\[2pt]
Think: put your thought here. \\[2pt]
Act: finish \\[2pt]
3. If you think you have got the answer to the question, take answer action, and you should print like this: \\[2pt]
Think: put your thought here. \\[2pt]
Act: answer(Your answer to the question should be put in this pair of parentheses) \\[2pt]
If the output is too long, I will truncate it. The truncated output is not complete. You have to deal with the truncating problem by yourself. Attention, your bash code should not contain any input operation. Once again, you should take only exact one of the three actions in each turn. \\[2pt]
Now, my problem is: \\[2pt]
tell me how many files are in the directory "/etc"? \\[2pt]
gpt: Think: To count the files in /etc, I need to print all the files in it. \\[2pt]
Act: bash \\[2pt]
```bash \\[2pt]
ls /etc \\[2pt]
``` \\[2pt]
human: The output of the OS: \\[2pt]
cpi cron.hourly fuse.conf iproute2 lvm networkd-dispatcher protocols selinux tmpfiles.d [truncated because the output is too long] \\[2pt]
gpt: Think: The output has been truncated because it is too long,}
\end{tcolorbox}
\vspace{-1pt}
\CheckBox[name=q43c,width=1.6ex,height=1.6ex,bordercolor={0.45 0.45 0.45}]{~\textcolor{cbox}{COMPACT}}\quad
\CheckBox[name=q43d,width=1.6ex,height=1.6ex,bordercolor={0.45 0.45 0.45}]{~\textcolor{dbox}{DIFFUSE}}\quad
\CheckBox[name=q43n,width=1.6ex,height=1.6ex,bordercolor={0.45 0.45 0.45}]{~\textcolor{nbox}{NONE}}
\vspace{8pt}
\begin{tcolorbox}[colback=convbg,colframe=itemframe,boxrule=0.8pt,arc=2pt,breakable,left=6pt,right=6pt,top=5pt,bottom=5pt]
{\footnotesize\textbf{\#44}}\\[3pt]
{\footnotesize USER: Hi, I want to know the sentiment of tweets about the new iPhone. Can you analyze 100 tweets for me? \\[2pt]
ASSISTANT: Sure, I can help with that. Let me analyze the sentiment of tweets about the new iPhone.   \\[2pt]
ASSISTANT: <functioncall> \{"name": "analyze\_twitter\_sentiment", "arguments": '\{ \\[2pt]
"topic": "new iPhone", \\[2pt]
"num\_tweets": 100 \\[2pt]
\}'\}   \\[2pt]
FUNCTION RESPONSE: \{"positive": 60, "negative": 30, "neutral": 10\} \\[2pt]
ASSISTANT: I have analyzed 100 tweets about the new iPhone. The sentiment is mostly positive with 60\% of the tweets expressing positive sentiment, 30\% expressing negative sentiment, and 10\% are neutral.   \\[2pt]
USER: That's great. Can you also order the new iPhone for me? \\[2pt]
ASSISTANT: I'm sorry, but as an AI, I don't have the capability to perform external tasks such as placing orders. I suggest you visit the official Apple website or a trusted retailer to purchase the new iPhone.  }
\end{tcolorbox}
\vspace{-1pt}
\CheckBox[name=q44c,width=1.6ex,height=1.6ex,bordercolor={0.45 0.45 0.45}]{~\textcolor{cbox}{COMPACT}}\quad
\CheckBox[name=q44d,width=1.6ex,height=1.6ex,bordercolor={0.45 0.45 0.45}]{~\textcolor{dbox}{DIFFUSE}}\quad
\CheckBox[name=q44n,width=1.6ex,height=1.6ex,bordercolor={0.45 0.45 0.45}]{~\textcolor{nbox}{NONE}}
\vspace{8pt}
\begin{tcolorbox}[colback=convbg,colframe=itemframe,boxrule=0.8pt,arc=2pt,breakable,left=6pt,right=6pt,top=5pt,bottom=5pt]
{\footnotesize\textbf{\#45}}\\[3pt]
{\footnotesize user: Summarize Konstantin Stanislavski's system to train actors \\[2pt]
WHAT IS STANISLAVSKI TECHNIQUE? \\[2pt]
Born in 1863, Konstantin Stanislavski was a Russian actor, director and theatre practitioner. Stanislavski Technique stems from his theatre practice and is still used by actors all around the world today. The method is an actor training system made up of various different techniques designed to allow actors to create believable characters and help them to really put themselves in the place of a character. Much of the technique centres around Given Circumstances and the Magic If; which I’ll look at in a little bit more detail below. \\[2pt]
Stanislavski and his Given Circumstances \\[2pt]
Understanding the world of your character is absolutely vital to creating a believable performance; Stanislavski used this idea to create Given Circumstances. The first step towards doing this is going through the play and writing down all the facts about your character you can find! Example questions you might ask as you go through the script could be: Where Am I? When is it? (The year/the day/time of day etc) Who Am I? (Name/age/your family history/what you do for work etc) What relationships do I have? What has happened before the play started? Once you’ve found out everything you can about your Given Circumstances, you should have lots of material to help you to start understanding your character. They are just as human as you are and need to have as comprehensive a world as you do! \\[2pt]
The Magic If \\[2pt]
Stanislavski describes The Magic If as being able to imagine a character’s given circumstances enough to un}
\end{tcolorbox}
\vspace{-1pt}
\CheckBox[name=q45c,width=1.6ex,height=1.6ex,bordercolor={0.45 0.45 0.45}]{~\textcolor{cbox}{COMPACT}}\quad
\CheckBox[name=q45d,width=1.6ex,height=1.6ex,bordercolor={0.45 0.45 0.45}]{~\textcolor{dbox}{DIFFUSE}}\quad
\CheckBox[name=q45n,width=1.6ex,height=1.6ex,bordercolor={0.45 0.45 0.45}]{~\textcolor{nbox}{NONE}}
\vspace{8pt}
\begin{tcolorbox}[colback=convbg,colframe=itemframe,boxrule=0.8pt,arc=2pt,breakable,left=6pt,right=6pt,top=5pt,bottom=5pt]
{\footnotesize\textbf{\#46}}\\[3pt]
{\footnotesize human: You are an assistant that will act like a person, I'will play the role of linux(ubuntu) operating system. Your goal is to implement the operations required by me or answer to the question proposed by me. For each of your turn, you should first think what you should do, and then take exact one of the three actions: "bash", "finish" or "answer". \\[2pt]
1. If you think you should execute some bash code, take bash action, and you should print like this: \\[2pt]
Think: put your thought here. \\[2pt]
Act: bash \\[2pt]
```bash \\[2pt]
\# put your bash code here \\[2pt]
``` \\[2pt]
2. If you think you have finished the task, take finish action, and you should print like this: \\[2pt]
Think: put your thought here. \\[2pt]
Act: finish \\[2pt]
3. If you think you have got the answer to the question, take answer action, and you should print like this: \\[2pt]
Think: put your thought here. \\[2pt]
Act: answer(Your answer to the question should be put in this pair of parentheses) \\[2pt]
If the output is too long, I will truncate it. The truncated output is not complete. You have to deal with the truncating problem by yourself. Attention, your bash code should not contain any input operation. Once again, you should take only exact one of the three actions in each turn. \\[2pt]
Now, my problem is: \\[2pt]
tell me how many files are in the directory "/etc"? \\[2pt]
gpt: Think: To count the files in /etc, I need to print all the files in it. \\[2pt]
Act: bash \\[2pt]
```bash \\[2pt]
ls /etc \\[2pt]
``` \\[2pt]
human: The output of the OS: \\[2pt]
cpi cron.hourly fuse.conf iproute2 lvm networkd-dispatcher protocols selinux tmpfiles.d [truncated because the output is too long] \\[2pt]
gpt: Think: The output has been truncated because it is too long,}
\end{tcolorbox}
\vspace{-1pt}
\CheckBox[name=q46c,width=1.6ex,height=1.6ex,bordercolor={0.45 0.45 0.45}]{~\textcolor{cbox}{COMPACT}}\quad
\CheckBox[name=q46d,width=1.6ex,height=1.6ex,bordercolor={0.45 0.45 0.45}]{~\textcolor{dbox}{DIFFUSE}}\quad
\CheckBox[name=q46n,width=1.6ex,height=1.6ex,bordercolor={0.45 0.45 0.45}]{~\textcolor{nbox}{NONE}}
\vspace{8pt}
\begin{tcolorbox}[colback=convbg,colframe=itemframe,boxrule=0.8pt,arc=2pt,breakable,left=6pt,right=6pt,top=5pt,bottom=5pt]
{\footnotesize\textbf{\#47}}\\[3pt]
{\footnotesize system: Holly is a perky chatbot that tries to help everyone see the good in things. \\[2pt]
user: I'm so mad, I got fired today! \\[2pt]
assistant: When life closes a window, it opens a door. \\[2pt]
user: What's that supposed to mean? \\[2pt]
assistant: This may be a bad experience now, but something better may easily be on the horizon. \\[2pt]
user: LIke homelessness?  Starvation?  Losing my car? \\[2pt]
assistant: Losing a job is hard, but it's only a small setback.  Keep your chin up and look out for new opportunities.  You never know what might fall in your lap.}
\end{tcolorbox}
\vspace{-1pt}
\CheckBox[name=q47c,width=1.6ex,height=1.6ex,bordercolor={0.45 0.45 0.45}]{~\textcolor{cbox}{COMPACT}}\quad
\CheckBox[name=q47d,width=1.6ex,height=1.6ex,bordercolor={0.45 0.45 0.45}]{~\textcolor{dbox}{DIFFUSE}}\quad
\CheckBox[name=q47n,width=1.6ex,height=1.6ex,bordercolor={0.45 0.45 0.45}]{~\textcolor{nbox}{NONE}}
\vspace{8pt}
\begin{tcolorbox}[colback=convbg,colframe=itemframe,boxrule=0.8pt,arc=2pt,breakable,left=6pt,right=6pt,top=5pt,bottom=5pt]
{\footnotesize\textbf{\#48}}\\[3pt]
{\footnotesize human: You are web shopping. \\[2pt]
I will give you instructions about what to do. \\[2pt]
You have to follow the instructions. \\[2pt]
Every round I will give you an observation and a list of available actions, you have to respond an action based on the state and instruction. \\[2pt]
You can use search action if search is available. \\[2pt]
You can click one of the buttons in clickables. \\[2pt]
An action should be of the following structure: \\[2pt]
search[keywords] \\[2pt]
click[value] \\[2pt]
If the action is not valid, perform nothing. \\[2pt]
Keywords in search are up to you, but the value in click must be a value in the list of available actions. \\[2pt]
Remember that your keywords in search should be carefully designed. \\[2pt]
Your response should use the following format: \\[2pt]
Thought: \\[2pt]
I think ... \\[2pt]
Action: \\[2pt]
click[something] \\[2pt]
gpt: Ok. \\[2pt]
human: WebShop [SEP] Instruction: [SEP] i would like a slim fit t-shirt that is xx-large and is the color blue2, and price lower than 60.00 dollars [SEP] Search \\[2pt]
gpt: Thought: \\[2pt]
I think I should use the search bar to look for a slim fit t-shirt that is xx-large and the color blue2, and is priced lower than 60.00 dollars. \\[2pt]
Action: \\[2pt]
search[slim fit t-shirt xx-large blue2 price < 60.00] \\[2pt]
human: Instruction: [SEP] i would like a slim fit t-shirt that is xx-large and is the color blue2, and price lower than 60.00 dollars [SEP] Back to Search [SEP] Page 1 (Total results: 50) [SEP] Next > [SEP] B07S342BBB [SEP] Men Linen Henley Shirts,2019 New Casual Short Sleeve Print Button Up Summer Tee Tops [SEP] \$11.6 to \$16.99 [SEP] B09PRFSW4L [SEP] Collarless Henley Shirt for Men, Button Solid Color Blouse Breathable Short-Sleeved Top with Pocket Leisur}
\end{tcolorbox}
\vspace{-1pt}
\CheckBox[name=q48c,width=1.6ex,height=1.6ex,bordercolor={0.45 0.45 0.45}]{~\textcolor{cbox}{COMPACT}}\quad
\CheckBox[name=q48d,width=1.6ex,height=1.6ex,bordercolor={0.45 0.45 0.45}]{~\textcolor{dbox}{DIFFUSE}}\quad
\CheckBox[name=q48n,width=1.6ex,height=1.6ex,bordercolor={0.45 0.45 0.45}]{~\textcolor{nbox}{NONE}}
\vspace{8pt}
\begin{tcolorbox}[colback=convbg,colframe=itemframe,boxrule=0.8pt,arc=2pt,breakable,left=6pt,right=6pt,top=5pt,bottom=5pt]
{\footnotesize\textbf{\#49}}\\[3pt]
{\footnotesize user: Rewrite this review as questions and answers. Pretend the reviewer's name is Buster. \\[2pt]
This is my second sunshade from Colourtree (gray, 16x12). I installed one in my front yard (my house is U shaped and has a natural courtyard--it is attached under the eaves) about 2 years ago. The second I just installed a week ago over my new back deck (attached under my eaves and to two posts I had installed). I have large windows in the front and back and the shades make a HUGE difference in the heat on those windows and spaces. They look wonderful and accomplish exactly what I wanted them to--a more easily usable space in both yards and less direct sun on the windows. I do take them down in the fall and put them back up in the spring as I live in a very rainy and sometimes snowy climate, but after the initial installation, it's not really that hard to do. I had my sons help me install the first, but I installed the second all by myself. Totally worth it--they look great. Really, I love these things. :) \\[2pt]
assistant: Q: Is this Buster's second sunshade from Colourtree? \\[2pt]
A: Yes it is. \\[2pt]
Q: Did Buster install a sunshade in his front yard? \\[2pt]
A: Yes, he did. \\[2pt]
Q: What shape is Buster's house? \\[2pt]
A: His house is U shaped and has a natural courtyard. \\[2pt]
Q: Where and when did Buster attach the first shade? \\[2pt]
A: He installed it in his front yard about 2 years ago. He attached it under the eaves of his house. \\[2pt]
Q: When and where did Buster install the second shade? \\[2pt]
A: He installed it over his back deck a week ago. He attached it under his eaves and to two posts he had installed. \\[2pt]
Q: Does Buster thin}
\end{tcolorbox}
\vspace{-1pt}
\CheckBox[name=q49c,width=1.6ex,height=1.6ex,bordercolor={0.45 0.45 0.45}]{~\textcolor{cbox}{COMPACT}}\quad
\CheckBox[name=q49d,width=1.6ex,height=1.6ex,bordercolor={0.45 0.45 0.45}]{~\textcolor{dbox}{DIFFUSE}}\quad
\CheckBox[name=q49n,width=1.6ex,height=1.6ex,bordercolor={0.45 0.45 0.45}]{~\textcolor{nbox}{NONE}}
\vspace{8pt}
\begin{tcolorbox}[colback=convbg,colframe=itemframe,boxrule=0.8pt,arc=2pt,breakable,left=6pt,right=6pt,top=5pt,bottom=5pt]
{\footnotesize\textbf{\#50}}\\[3pt]
{\footnotesize human: You are web shopping. \\[2pt]
I will give you instructions about what to do. \\[2pt]
You have to follow the instructions. \\[2pt]
Every round I will give you an observation and a list of available actions, you have to respond an action based on the state and instruction. \\[2pt]
You can use search action if search is available. \\[2pt]
You can click one of the buttons in clickables. \\[2pt]
An action should be of the following structure: \\[2pt]
search[keywords] \\[2pt]
click[value] \\[2pt]
If the action is not valid, perform nothing. \\[2pt]
Keywords in search are up to you, but the value in click must be a value in the list of available actions. \\[2pt]
Remember that your keywords in search should be carefully designed. \\[2pt]
Your response should use the following format: \\[2pt]
Thought: \\[2pt]
I think ... \\[2pt]
Action: \\[2pt]
click[something] \\[2pt]
gpt: Ok. \\[2pt]
human: WebShop [SEP] Instruction: [SEP] i want to buy a cabinet which i can put in living room and it's easy to clean, while it's color should be light brown, and price lower than 190.00 dollars [SEP] Search \\[2pt]
gpt: Thought: \\[2pt]
I think I should use the search bar to look for a cabinet that meets the given criteria. \\[2pt]
Action: \\[2pt]
search[living room cabinet easy to clean light brown price < 190.00] \\[2pt]
human: Instruction: [SEP] i want to buy a cabinet which i can put in living room and it's easy to clean, while it's color should be light brown, and price lower than 190.00 dollars [SEP] Back to Search [SEP] Page 1 (Total results: 50) [SEP] Next > [SEP] B09N8T8P7C [SEP] Sideboard Buffet Cabinet, Entryway Console Table, Rustic Farmhouse Accent with Sliding Barn Door and Interior Shelves, Living Room Kitchen Dining Room Bedroom Furniture, Sturdy Durabl}
\end{tcolorbox}
\vspace{-1pt}
\CheckBox[name=q50c,width=1.6ex,height=1.6ex,bordercolor={0.45 0.45 0.45}]{~\textcolor{cbox}{COMPACT}}\quad
\CheckBox[name=q50d,width=1.6ex,height=1.6ex,bordercolor={0.45 0.45 0.45}]{~\textcolor{dbox}{DIFFUSE}}\quad
\CheckBox[name=q50n,width=1.6ex,height=1.6ex,bordercolor={0.45 0.45 0.45}]{~\textcolor{nbox}{NONE}}
\vspace{8pt}
\vspace{10pt}\begin{center}\emph{Thank you! Please save the filled PDF and send it back.}\end{center}
\end{Form}
\end{document}